% Standalone LaTeX edition generated from gpu_systems_crash_course.md.
% The H100 comparison table uses explicit wrapping columns for print layout.
% Options for packages loaded elsewhere
\PassOptionsToPackage{unicode}{hyperref}
\PassOptionsToPackage{hyphens}{url}
\PassOptionsToPackage{dvipsnames,svgnames,x11names}{xcolor}
\documentclass[
  10pt,
  letterpaper,
]{article}
\usepackage{xcolor}
\usepackage[margin=0.8in]{geometry}
\usepackage{amsmath,amssymb}
\setcounter{secnumdepth}{-\maxdimen} % remove section numbering
\usepackage{iftex}
\ifPDFTeX
  \usepackage[T1]{fontenc}
  \usepackage[utf8]{inputenc}
  \usepackage{textcomp} % provide euro and other symbols
\else % if luatex or xetex
  \usepackage{unicode-math} % this also loads fontspec
  \defaultfontfeatures{Scale=MatchLowercase}
  \defaultfontfeatures[\rmfamily]{Ligatures=TeX,Scale=1}
\fi
\usepackage{lmodern}
\ifPDFTeX\else
  % xetex/luatex font selection
\fi
% Use upquote if available, for straight quotes in verbatim environments
\IfFileExists{upquote.sty}{\usepackage{upquote}}{}
\IfFileExists{microtype.sty}{% use microtype if available
  \usepackage[]{microtype}
  \UseMicrotypeSet[protrusion]{basicmath} % disable protrusion for tt fonts
}{}
\makeatletter
\@ifundefined{KOMAClassName}{% if non-KOMA class
  \IfFileExists{parskip.sty}{%
    \usepackage{parskip}
  }{% else
    \setlength{\parindent}{0pt}
    \setlength{\parskip}{6pt plus 2pt minus 1pt}}
}{% if KOMA class
  \KOMAoptions{parskip=half}}
\makeatother
\usepackage{longtable,booktabs,array}
\newcounter{none} % for unnumbered tables
\usepackage{calc} % for calculating minipage widths
% Correct order of tables after \paragraph or \subparagraph
\usepackage{etoolbox}
\makeatletter
\patchcmd\longtable{\par}{\if@noskipsec\mbox{}\fi\par}{}{}
\makeatother
% Allow footnotes in longtable head/foot
\IfFileExists{footnotehyper.sty}{\usepackage{footnotehyper}}{\usepackage{footnote}}
\makesavenoteenv{longtable}
\setlength{\emergencystretch}{3em} % prevent overfull lines
\providecommand{\tightlist}{%
  \setlength{\itemsep}{0pt}\setlength{\parskip}{0pt}}
\usepackage{bookmark}
\IfFileExists{xurl.sty}{\usepackage{xurl}}{} % add URL line breaks if available
\urlstyle{same}
\hypersetup{
  pdftitle={GPU Systems Engineering: A Pareto-Oriented Crash Course and the PBG Case Study},
  colorlinks=true,
  linkcolor={blue},
  filecolor={Maroon},
  citecolor={Blue},
  urlcolor={blue},
  pdfcreator={LaTeX via pandoc}}

\title{GPU Systems Engineering: A Pareto-Oriented Crash Course and the
PBG Case Study}
\usepackage{etoolbox}
\makeatletter
\providecommand{\subtitle}[1]{% add subtitle to \maketitle
  \apptocmd{\@title}{\par {\large #1 \par}}{}{}
}
\makeatother
\subtitle{From CUDA fundamentals to the HU NLHE H100/RTX 2080 Ti case
study}
\author{}
\date{August 2026}

\begin{document}
\maketitle

{
\setcounter{tocdepth}{2}
\tableofcontents
}
Status: explanatory background and engineering history, not a
mathematical or runtime contract. The authoritative project contracts
remain the source code, \texttt{AGENTS.md}, \texttt{CONTRACTS.md}, and
the focused documents under \texttt{docs/queued/}. Hardware peak numbers
are specification ceilings, not performance predictions.

\subsection{Who this is for}\label{who-this-is-for}

This chapter assumes mature mathematical and computer-science instincts
but no prior systems, GPU, or CUDA experience. Familiarity with
asymptotic analysis, arrays, probability, and neural networks is
helpful. No knowledge of computer architecture beyond ``a CPU executes
instructions and memory stores bytes'' is required.

The first part is deliberately generic. It develops a mental model for
GPU systems and a method for finding high-return optimizations without
getting lost in low-level cleverness. The final part applies that model
to this repository\textquotesingle s NLHE learning pipeline, its RTX
2080 Ti and H100 environments, the failures the H100 work exposed, the
fixes that followed, and the experiments that remain worth considering.

Two uses of the letter \(R\) would otherwise collide. In the generic and
project hardware discussion, \(K\) denotes a number of independent model
replicas. In the poker configuration, \(R\) denotes the number of
retained card ranks.

\begin{center}\rule{0.5\linewidth}{0.5pt}\end{center}

\section{Part I --- A GPU and CUDA systems crash
course}\label{part-i--a-gpu-and-cuda-systems-crash-course}

\subsection{1. The big picture: a GPU is a throughput machine inside a
larger
system}\label{1-the-big-picture-a-gpu-is-a-throughput-machine-inside-a-larger-system}

The shortest useful description of a modern accelerator is:

\begin{quote}
A GPU is a large collection of relatively simple arithmetic lanes, fed
by a deep memory hierarchy and controlled by a host CPU. It is
exceptionally fast when the program exposes enough similar, independent
work and exceptionally easy to starve when it does not.
\end{quote}

A conventional CPU spends substantial silicon and power making a few
threads fast. Its cores have sophisticated branch prediction, large
caches, speculative and out-of-order execution, and strong single-thread
latency. A GPU spends more of its budget on arithmetic throughput and
hardware support for keeping many threads in flight. It tolerates the
latency of one memory request by running other ready work rather than
making one thread complete as quickly as possible.

That difference changes the optimization objective:

\begin{itemize}
\tightlist
\item
  CPU question: ``How quickly can this dependency chain finish?''
\item
  GPU question: ``Can thousands of arithmetic lanes keep receiving
  useful work?''
\end{itemize}

This is a tendency, not a law. CPUs can exploit vector and thread
parallelism; GPUs can execute dependency-heavy code. The point is
economic: each device is best when the program matches what the hardware
bought with its transistor and power budget.

\subsubsection{1.1 Latency, throughput, and utilization are different
quantities}\label{11-latency-throughput-and-utilization-are-different-quantities}

Suppose one request takes 10 ms on a CPU and 15 ms on a GPU. If the CPU
handles one request at a time while the GPU handles 1,000 together, the
GPU has worse single-request latency but vastly better aggregate
throughput. Conversely, a GPU capable of an enormous peak operation rate
may show poor utilization if it receives one tiny request every
millisecond.

Three metrics should therefore never be conflated:

\begin{enumerate}
\def\labelenumi{\arabic{enumi}.}
\tightlist
\item
  \textbf{Latency:} elapsed time for one unit of work.
\item
  \textbf{Throughput:} useful units completed per second.
\item
  \textbf{Utilization:} how busy some hardware resource appears.
\end{enumerate}

High utilization is not automatically high useful throughput. A GPU can
be busy performing padding, recomputation, copies, or a mathematically
wrong objective. Low utilization is not automatically a defect either: a
memory-bound kernel can saturate memory bandwidth while most arithmetic
units are idle. The real objective must be stated in domain units such
as examples/s, tokens/s, simulated decisions/s, or time-to-quality.

\subsubsection{1.2 The whole machine
matters}\label{12-the-whole-machine-matters}

CUDA code runs inside a heterogeneous system:

\begin{verbatim}
storage -> CPU process -> host memory -> driver/runtime -> GPU memory
        -> GPU kernels -> synchronization -> results/checkpoints
\end{verbatim}

The GPU cannot compensate for an input pipeline that stalls, a Python
loop that submits work too slowly, repeated host/device transfers, an
allocator that fragments memory, or a distributed collective waiting on
another rank. ``GPU optimization'' is consequently system optimization.
The unit of analysis is an end-to-end critical path, not an isolated
multiplication.

\subsubsection{1.3 The Pareto principle for accelerator
work}\label{13-the-pareto-principle-for-accelerator-work}

In practice, a small set of choices often controls most performance:

\begin{itemize}
\tightlist
\item
  keep the working set on the accelerator;
\item
  expose enough parallel work through batching or another honest
  independent axis;
\item
  use optimized library kernels for dense linear algebra and attention;
\item
  choose a supported low-precision format while keeping sensitive state
  in appropriate precision;
\item
  eliminate unintended synchronization and excessive tiny launches;
\item
  bound memory by actual lifetimes rather than by the sum of all
  possible objects; and
\item
  measure the real workload before writing custom kernels.
\end{itemize}

The hard part is not memorizing optimization tricks. It is identifying
which of these constraints is active in the current program. The Pareto
approach is to buy the largest credible end-to-end gain per unit of
implementation risk.

One useful decision score is informal but disciplined:

\[
\text{ROI score}
\approx
\frac{\text{possible end-to-end gain}\times\text{confidence}}
{\text{engineering cost}\times\text{correctness risk}\times\text{maintenance cost}}.
\]

No one should pretend that the factors are known precisely. The formula
is a reminder that a 50\% microkernel opportunity covering 2\% of
runtime is worth at most about 1\% end to end, while a mundane batching
change covering 60\% of runtime can dominate months of kernel work.

\subsubsection{1.4 A running example: one million independent
rows}\label{14-a-running-example-one-million-independent-rows}

To make the abstractions concrete, imagine that we have one million
rows. Each row contains a state vector \(x_i\in\mathbb{R}^{128}\), and
we want to compute

\[
y_i=\operatorname{ReLU}(x_iW+b),
\qquad W\in\mathbb{R}^{128\times128}.
\]

The rows are independent and share \(W\). Several implementations are
possible:

\begin{enumerate}
\def\labelenumi{\arabic{enumi}.}
\tightlist
\item
  A Python loop calls the model one row at a time.
\item
  A CPU thread pool assigns row ranges to cores.
\item
  The GPU receives one row at a time.
\item
  The GPU receives a large batch and calls one matrix multiplication.
\end{enumerate}

All four compute the same mathematical map. They present radically
different physical work. Option 3 pays host dispatch and kernel startup
repeatedly while giving each launch almost no parallel slack. Option 4
turns the dominant work into a large GEMM, reuses \(W\), and amortizes
the launch. This is the basic reason ``put it on the GPU'' is incomplete
advice: the presentation of the work usually matters more than the
existence of a CUDA device.

Now change the problem so every row has its own matrix \(W_i\). The
shared-weight GEMM is no longer valid. We need batched/grouped GEMM or
another explicit independent-model strategy. The fastest-looking
transformation can be mathematically wrong if it accidentally shares
weights. Throughout this chapter, ``honest parallel axis'' means an axis
that exposes concurrency while preserving exactly this kind of semantic
distinction.

\subsection{2. CUDA\textquotesingle s execution model, from function
call to active
warps}\label{2-cudas-execution-model-from-function-call-to-active-warps}

CUDA is both a programming model and a software platform. The CPU is the
\textbf{host}; an NVIDIA GPU is the \textbf{device}. Host code allocates
memory, prepares arguments, launches device functions called
\textbf{kernels}, and coordinates dependencies. A kernel launch
conceptually creates a grid of lightweight threads.

\subsubsection{2.1 Grid, block, thread}\label{21-grid-block-thread}

CUDA threads are organized hierarchically:

\begin{verbatim}
kernel launch
  grid
    thread block 0
      thread 0, thread 1, ...
    thread block 1
      thread 0, thread 1, ...
    ...
\end{verbatim}

Threads in a block can cooperate through fast on-chip shared memory and
block-level barriers. Blocks should be independent enough that the
hardware may run them in any order on available \textbf{streaming
multiprocessors} (SMs). This independence is what lets the same kernel
scale across GPUs with different numbers of SMs.

A vector addition illustrates the indexing model:

\begin{verbatim}
__global__ void add(const float* x, const float* y, float* z, int n) {
    int i = blockIdx.x * blockDim.x + threadIdx.x;
    if (i < n) z[i] = x[i] + y[i];
}
\end{verbatim}

One launch creates enough blocks to cover \(n\) elements. Each logical
thread computes its global index and handles one element. Real kernels
often make each thread do more work, tile data cooperatively, or use
multiple dimensions, but the hierarchy remains.

\subsubsection{2.2 Warps and SIMT}\label{22-warps-and-simt}

The hardware schedules threads in groups of 32 called \textbf{warps}.
CUDA uses the term SIMT: single instruction, multiple threads. Threads
retain their own register state and may follow different logical paths,
but a warp is most efficient when its active threads execute the same
instruction on adjacent data.

If half a warp takes one branch and half takes another, the hardware
generally executes the necessary paths with different active masks. This
\textbf{divergence} does not make the result wrong, but it can reduce
useful work per issued instruction. Divergence between different warps
is unremarkable; divergence within a warp is the concern.

This suggests a general design pattern: group similar work together.
Sort or bucket rows by shape, phase, or operation type before a large
kernel when the cost of grouping is smaller than the waste it removes.
Do not sort merely for aesthetic regularity; measure the trade.

\subsubsection{2.3 SMs, resident blocks, and latency
hiding}\label{23-sms-resident-blocks-and-latency-hiding}

An SM owns execution pipelines, warp schedulers, a register file,
caches, and shared memory. Multiple blocks and warps may be resident on
an SM. When one warp waits for data or an instruction dependency, a
scheduler can issue a ready warp. This is \textbf{latency hiding}.

The fraction of the hardware\textquotesingle s theoretical resident-warp
capacity actually occupied is called \textbf{occupancy}. Occupancy is
limited by resources including:

\begin{itemize}
\tightlist
\item
  threads per block;
\item
  registers per thread;
\item
  shared memory per block; and
\item
  architectural block and warp limits.
\end{itemize}

Higher occupancy can provide more latency-hiding opportunities, but
100\% occupancy is not the objective. A kernel using more registers may
do less spilling and finish sooner at lower occupancy. A compute-dense
matrix kernel can perform well without maximal residency. Occupancy is a
diagnostic constraint, not a score.

\subsubsection{2.4 Ordinary cores and Tensor
Cores}\label{24-ordinary-cores-and-tensor-cores}

Modern NVIDIA GPUs contain general arithmetic pipelines and specialized
\textbf{Tensor Cores} for matrix multiply-accumulate operations on
supported shapes and data types. Neural networks spend much of their
time in matrix products, so Tensor Cores can provide much higher
throughput than scalar FP32 pipelines.

They are not magic neural-network cores. A framework or kernel must
present compatible operations, layouts, dimensions, alignments, and data
types. Small, awkward, or fragmented matrix products may realize only a
small fraction of the published peak. ``The GPU has 989 dense BF16
Tensor TFLOP/s'' and ``this model can perform 989 trillion useful
operations per second'' are radically different claims.

\subsubsection{2.5 Compute capability is an executable
contract}\label{25-compute-capability-is-an-executable-contract}

NVIDIA labels architectural feature sets with a \textbf{compute
capability}. Turing devices such as the RTX 2080 Ti use 7.5; H100 uses
9.0. This affects available instructions, supported numeric formats,
shared-memory limits, and which compiled device binaries can execute.

Compute capability is not the installed CUDA toolkit version. A machine
can have a new toolkit and an old GPU. Native CUDA extensions should be
compiled for their target architectures, and build caches should not
silently reuse an SM75 artifact for an SM90 experiment. PTX just-in-time
compilation can provide some forward compatibility, but it does not
eliminate the need to verify the actual artifact and dispatch path.

\subsubsection{2.6 Worked launch
geometry}\label{26-worked-launch-geometry}

Return to a simple elementwise kernel over \(n=1{,}000{,}000\) values.
With 256 threads per block, the grid needs

\[
\left\lceil\frac{1{,}000{,}000}{256}\right\rceil=3{,}907
\]

blocks. The last block contains inactive threads beyond the array
boundary; the \texttt{if\ (i\ \textless{}\ n)} guard makes them
harmless. The runtime distributes blocks to SMs as resources become
available. The source does not hard-code whether 68 or 132 SMs exist,
which is how the same grid scales across the RTX 2080 Ti and H100.

Why not launch one block of one million threads? A block has an
architectural thread limit and must reside on one SM. Thousands of
independent blocks let the hardware load-balance across every SM. Why
not use one thread per block? That would create enormous scheduling
overhead and give each warp only one useful lane. Blocks commonly use a
multiple of 32 threads so warps are full, but the best block size also
depends on registers, shared memory, and work per thread.

Suppose a kernel uses so many registers that only two blocks can reside
on an SM. Reducing block size might permit more blocks and warps, but
could also make each block less efficient. The practical method is to
inspect compiler resource usage, use an occupancy calculator as a
constraint check, and benchmark a small set of sensible launch
configurations. ``256 threads'' is a common starting point, not a
universal constant.

\subsection{3. The memory hierarchy: where performance is usually won or
lost}\label{3-the-memory-hierarchy-where-performance-is-usually-won-or-lost}

Arithmetic is useful only when operands arrive. A simplified device
hierarchy is:

\begin{verbatim}
registers (private to a thread; fastest and smallest)
    ↓
shared memory / L1 cache (on an SM)
    ↓
L2 cache (shared across SMs)
    ↓
device DRAM: GDDR6 or HBM (large, high bandwidth, high latency)
    ↓
host-device link: PCIe or a platform interconnect
    ↓
host DRAM and storage
\end{verbatim}

The arrows should not be interpreted as a mandatory access sequence;
caches and special paths complicate the machine. The important lesson is
that distance and capacity increase together, while latency and the
energy cost of moving a byte generally increase as well.

\subsubsection{3.1 Bandwidth, latency, and
capacity}\label{31-bandwidth-latency-and-capacity}

\begin{itemize}
\tightlist
\item
  \textbf{Capacity} asks whether the working set fits.
\item
  \textbf{Bandwidth} asks how many bytes per second can move under
  sustained load.
\item
  \textbf{Latency} asks how long one dependency takes.
\end{itemize}

An algorithm can fit in memory yet be bandwidth-bound. It can use little
bandwidth yet stall on dependent random loads. Adding HBM capacity
solves an out-of-memory problem but does not automatically improve a
launch-bound tiny model. These axes must be diagnosed separately.

\subsubsection{3.2 Coalescing and structure of
arrays}\label{32-coalescing-and-structure-of-arrays}

When threads in a warp access adjacent addresses, the hardware can
combine the requests into relatively few memory transactions. Scattered
or strided access can fetch many bytes that no thread uses. This
property is called \textbf{coalescing}.

Consider 32 poker states. An array-of-structures layout stores all
fields for state 0, then all fields for state 1:

\begin{verbatim}
[state0: actor, street, pot, ...][state1: actor, street, pot, ...]...
\end{verbatim}

If a kernel needs only \texttt{actor}, consecutive threads jump by the
full structure width. A structure-of-arrays layout stores one dense
array per field:

\begin{verbatim}
actor:  [a0, a1, a2, ...]
street: [s0, s1, s2, ...]
pot:    [p0, p1, p2, ...]
\end{verbatim}

Now adjacent threads read adjacent actors. SoA also permits fields to
have different types and lifetimes. It is not universally superior---one
thread using every field of one small object may prefer AoS or a
hybrid---but it is a powerful default for batched GPU pipelines.

\subsubsection{3.3 Reuse and arithmetic
intensity}\label{33-reuse-and-arithmetic-intensity}

Define arithmetic intensity as

\[
I=\frac{\text{arithmetic operations}}{\text{bytes transferred from the relevant memory level}}.
\]

A matrix multiplication reuses tiles of its inputs many times, so its
intensity can be high. An elementwise addition reads two values and
writes one for one addition, so its intensity is low. Fusion can raise
intensity by keeping an intermediate in registers or shared memory
instead of writing it to DRAM and reading it back.

The qualifier ``relevant memory level'' matters. A kernel may be
HBM-efficient but limited by L2 traffic, shared-memory bandwidth, or
register dependencies. The basic roofline model introduced later usually
starts with device DRAM and then becomes more detailed when evidence
requires it.

\subsubsection{3.4 Allocated, reserved, live, and peak
memory}\label{34-allocated-reserved-live-and-peak-memory}

Framework memory reports are easy to misread:

\begin{itemize}
\tightlist
\item
  \textbf{Live/allocated tensor memory:} bytes currently owned by live
  tensors.
\item
  \textbf{Reserved memory:} blocks retained by a caching allocator for
  reuse.
\item
  \textbf{Peak allocated/reserved:} high-water marks since counters were
  reset.
\item
  \textbf{Workspace:} temporary library or kernel storage, sometimes
  hidden behind an operator boundary.
\item
  \textbf{Fragmentation/non-releasable memory:} free space inside
  reserved blocks that cannot satisfy the next request or be returned
  conveniently.
\end{itemize}

Reserved memory remaining after a step is not by itself a leak.
Monotonic growth across equivalent steady-state iterations is more
suspicious. A first iteration may also understate the eventual peak if
later iterations create a reference cache, compiled graph pool,
optimizer state, or new shape class.

\subsubsection{3.5 Lifetime analysis beats ``sum every object''
budgeting}\label{35-lifetime-analysis-beats-sum-every-object-budgeting}

Suppose rollout needs 20 GB, evaluator completion needs 30 GB, and
learning needs 40 GB. If the phases are serialized and earlier
phase-only storage is released before the next phase, the lower bound is
near the largest compatible phase envelope, not \(20+30+40=90\) GB. On
the other hand, a tensor retained for backward or checkpoint rollback
remains live even if the source-level variable name goes out of scope.

A useful memory design starts with a lifetime table:

{\def\LTcaptype{none} % do not increment counter
\begin{longtable}[]{@{}lllrl@{}}
\toprule\noalign{}
Object & Created & Last consumer & Replicated? & Size law \\
\midrule\noalign{}
\endhead
\bottomrule\noalign{}
\endlastfoot
model parameters & session start & session end & \(K\) & \(O(KN)\) \\
rollout state & rollout & sealing & \(K\) & workload-dependent \\
completion workspace & completion & completion publication & often \(K\)
& samples-dependent \\
activations & learner forward & backward & active cohort &
\(O(KBNL)\) \\
optimizer moments & first update & session end & \(K\) & \(O(KN)\) \\
\end{longtable}
}

The last-consumer column is more useful than the creation site. It
determines when memory can actually be reused.

\subsubsection{3.6 Worked memory
examples}\label{36-worked-memory-examples}

Consider a warp whose 32 threads each load one adjacent FP32 value. The
warp requests 128 useful bytes. On devices of compute capability 6.0 or
newer, those loads can be serviced by four aligned 32-byte transactions.
If instead thread \(t\) loads element \(32t\), the warp touches 32
widely separated regions and may cause roughly one transaction per
useful four-byte value. The arithmetic is identical; the second layout
moves far more memory.

Now consider the elementwise operation

\[
y_i=ax_i+y_i.
\]

Ignoring cache reuse, it performs two floating operations and moves
about 12 bytes per FP32 element: read \(x_i\), read \(y_i\), write
\(y_i\). Its arithmetic intensity is only

\[
I\approx\frac{2}{12}=0.167\ \text{FLOP/byte}.
\]

At the H100 SXM specification bandwidth of 3.35 TB/s, even perfect
streaming would give a device-memory roof around \(0.56\) TFLOP/s for
this operation---far below the GPU\textquotesingle s FP32 or Tensor Core
peak. It is not a bad kernel merely because arithmetic utilization is
low; it is a memory operation.

Finally, suppose three consecutive pointwise kernels each read and write
a one-gigabyte tensor. In the worst simple accounting they cause several
gigabytes of avoidable round trips. Fusing them can be worth more than
tuning the arithmetic instructions because the saved bytes dominate.
This is why a profiler needs both kernel time and traffic, not only
FLOP/s.

\subsection{4. A quantitative performance
model}\label{4-a-quantitative-performance-model}

Optimization is much easier when ``slow'' is decomposed. For an
end-to-end iteration, write schematically

\[
T_{\rm total}
=T_{\rm host}+T_{\rm launch}+T_{\rm compute}+T_{\rm memory}
+T_{\rm transfer}+T_{\rm sync}+T_{\rm communication}+T_{\rm I/O},
\]

with the warning that terms can overlap and therefore do not always add
as independent wall times. A trace should reveal the critical path and
the amount of overlap.

\subsubsection{4.1 Amdahl\textquotesingle s law: optimize coverage, not
glamour}\label{41-amdahls-law-optimize-coverage-not-glamour}

If fraction \(p\) of runtime is improved by factor \(s\), the ideal
total speedup is

\[
S=\frac{1}{(1-p)+p/s}.
\]

Even making 5\% of runtime infinitely fast yields only
\(1/(0.95)\approx1.053\times\). Making a 60\% component twice as fast
yields \(1/(0.4+0.3)\approx1.43\times\). Before implementing a custom
kernel, measure its coverage and calculate its maximum possible
end-to-end value.

Amdahl\textquotesingle s law describes fixed-work or strong-scaling
intuition. Gustafson\textquotesingle s law captures another important
possibility: a larger accelerator may let us solve a proportionally
larger problem in similar time. In simulation and learning, that can be
the more valuable use---more independent trajectories, hands, or
experiments---provided the extra work is scientifically useful.

\subsubsection{4.2 The roofline model}\label{42-the-roofline-model}

Let \(P_{\max}\) be peak arithmetic throughput and \(B_{\max}\)
sustained memory bandwidth. For arithmetic intensity \(I\), a simple
upper bound is

\[
P_{\rm attainable}\leq\min(P_{\max}, I B_{\max}).
\]

At low intensity, the sloped \(IB_{\max}\) roof indicates memory
pressure. At high intensity, the flat \(P_{\max}\) roof indicates
compute pressure. The intersection

\[
I^*=P_{\max}/B_{\max}
\]

is the ridge point.

For the H100 SXM dense BF16 specification, the rough HBM ridge point is

\[
I^*\approx\frac{989\ \text{TFLOP/s}}{3.35\ \text{TB/s}}
\approx295\ \text{FLOP/byte}.
\]

This uses two specification ceilings and is therefore only a first-order
guide, but it conveys scale. The preceding SAXPY-like operation at
\(0.167\) FLOP/byte lies overwhelmingly on the bandwidth side. A large
tiled matrix multiplication may reuse inputs enough to exceed hundreds
of FLOPs per HBM byte and approach the compute side. A small or skinny
GEMM may have adequate algorithmic intensity but still miss peak because
its shape cannot fill Tensor Core tiles or all SMs.

This model is valuable because it prevents category mistakes. If a
kernel moves 12 bytes for one FP32 addition, buying more Tensor Core
throughput cannot help much. If a large matrix multiplication is
compute-bound, shaving a small metadata read is unlikely to matter. Real
machines have several roofs---L1, L2, HBM, scalar pipelines, Tensor
Cores, instruction issue---and achieved ceilings depend on shape, but
the first-order classification is already useful.

\subsubsection{4.3 Launch-bound work}\label{43-launch-bound-work}

A third regime is especially important for small models: the GPU may be
neither compute-bound nor bandwidth-bound because kernels are too short
and the host cannot submit them without gaps. If a kernel runs for a few
microseconds and Python/framework dispatch consumes a comparable
interval, a timeline shows thin GPU rectangles separated by empty space.

Possible responses include:

\begin{itemize}
\tightlist
\item
  increase useful batch size;
\item
  fuse adjacent pointwise operations;
\item
  compile a stable region;
\item
  use a vectorized or \texttt{foreach} optimizer;
\item
  capture a repeated static schedule in a CUDA Graph; or
\item
  move a host-controlled numerical loop into one device operation.
\end{itemize}

The order matters. CUDA Graphs can replay a bad collection of tiny,
redundant kernels with lower launch cost, but they do not remove the
redundant work.

\subsubsection{4.4 Parallel slack and the saturation
curve}\label{44-parallel-slack-and-the-saturation-curve}

Let \(W\) be useful work submitted together. Throughput commonly follows
a curve that rises quickly and then plateaus:

\begin{verbatim}
throughput
    |                 __________
    |             ___/
    |          __/
    |       __/
    |______/_____________________ useful work in flight
\end{verbatim}

Below saturation, a larger batch, more roots, more sequences, or more
replicas can amortize fixed overhead and fill more SMs. Near saturation,
extra work adds memory pressure and latency without much throughput.
Beyond capacity, it causes out-of-memory failure or forces
microbatching.

The experimenter\textquotesingle s job is to find the plateau in domain
throughput, not to maximize a batch-size knob.

\subsubsection{4.5 Queueing intuition}\label{45-queueing-intuition}

If work arrives at average rate \(\lambda\) and spends average time
\(W\) in a stable subsystem, Little\textquotesingle s law gives average
in-flight work

\[
L=\lambda W.
\]

This is a useful systems sanity check. A pipeline seeking high
throughput must keep enough work in flight to cover service latency. But
queues that grow without bound indicate arrival rate above service rate;
they increase memory and latency rather than throughput. Fixed-capacity
queues and backpressure make this failure explicit.

\subsection{5. The high-ROI optimization
ladder}\label{5-the-high-roi-optimization-ladder}

The following order is not absolute, but it is a strong default because
each stage improves the evidence needed for the next.

\subsubsection{5.1 Level 0: define correct useful
work}\label{51-level-0-define-correct-useful-work}

Before making anything faster, define what must remain invariant:

\begin{itemize}
\tightlist
\item
  objective and normalization;
\item
  sampling distribution and RNG identity requirements;
\item
  number of optimizer transitions;
\item
  data and model versions;
\item
  failure atomicity and checkpoint semantics;
\item
  accepted numerical tolerance; and
\item
  domain work counters.
\end{itemize}

A faster biased sampler or half-sized training schedule is a different
experiment, not a systems optimization. Padding may be physically
executed but must not be counted as useful presentations. Recomputed
validation passes must not be hidden in a throughput denominator.

The best performance regression test is often an identity such as

\[
\text{presentations}=K\times\text{steps}\times\text{batch size},
\]

plus exact per-phase work ledgers and model-state comparisons.

\subsubsection{5.2 Level 1: measure end to end and by
phase}\label{52-level-1-measure-end-to-end-and-by-phase}

Record synchronized wall time around the user-visible operation and
device event time around GPU ranges. Break the iteration into phases
that correspond to real ownership boundaries. Count useful work, not
inferred work.

At minimum, collect:

\begin{itemize}
\tightlist
\item
  wall time and GPU time;
\item
  domain throughput;
\item
  allocated/reserved current and peak memory;
\item
  phase durations;
\item
  batch and shape histograms;
\item
  host/device copy bytes;
\item
  synchronization count;
\item
  kernel count and visible launch gaps; and
\item
  hardware/software identity.
\end{itemize}

Warm compilation, autotuning, lazy initialization, and allocator state
before steady-state timing. Report cold-start costs separately because
they are real but answer a different question.

\subsubsection{5.3 Level 2: remove movement and synchronization
mistakes}\label{53-level-2-remove-movement-and-synchronization-mistakes}

High-confidence wins include:

\begin{itemize}
\tightlist
\item
  keep intermediate numerical state on the device;
\item
  batch small host/device transfers;
\item
  avoid \texttt{.item()}, \texttt{.cpu()}, or printing a device tensor
  inside a hot loop;
\item
  replace device-wide synchronization with stream dependencies or a
  later necessary boundary;
\item
  use contiguous/coalesced layouts for dominant access patterns; and
\item
  allocate stable arenas or reuse buffers rather than repeatedly
  allocating variable temporary objects.
\end{itemize}

Some host reads are necessary for orchestration, error handling, or
reports. The goal is to place them at explicit boundaries rather than
let innocent-looking framework calls serialize every step.

\subsubsection{5.4 Level 3: expose more useful parallel
work}\label{54-level-3-expose-more-useful-parallel-work}

This is often the largest win. Candidate axes include examples, tokens,
trajectories, simulation roots, parameter replicas, environments, or
independent requests. Choose an axis whose aggregation does not change
the mathematics.

There is a crucial difference between \textbf{logical batch size} and
\textbf{physical microbatch size}. A logical batch defines the
statistical update. Physical microbatches are an execution decomposition
used to fit memory. If a logical batch of 256 is split into physical
slices, loss weights must reproduce the same update; otherwise memory
tuning has silently changed the objective.

\subsubsection{5.5 Level 4: use the best existing
primitives}\label{55-level-4-use-the-best-existing-primitives}

Prefer mature libraries for GEMM, batched GEMM, reductions, convolution,
and scaled dot-product attention. They encode years of
architecture-specific tuning. Verify the actual dispatch rather than
assuming an API name guarantees a particular kernel.

For example, PyTorch scaled dot-product attention can select Flash,
memory-efficient, or math implementations depending on dtype, shape,
mask, layout, device, and version. Configuring \texttt{sdpa} establishes
the requested operator, not proof that every call used Flash attention.

\subsubsection{5.6 Level 5: reduce framework and launch
overhead}\label{56-level-5-reduce-framework-and-launch-overhead}

Once work and shapes are stable:

\begin{itemize}
\tightlist
\item
  use \texttt{torch.compile} on bounded hot regions;
\item
  inspect graph breaks and recompilations;
\item
  use \texttt{foreach} or fused optimizers where semantics match;
\item
  capture repeated static schedules in CUDA Graphs; and
\item
  cache only a bounded number of shape-specialized variants.
\end{itemize}

Compilation can fuse elementwise work and reduce Python overhead. CUDA
Graphs move repeated launch construction out of the replay path. Both
can increase memory use through cached code, workspaces, and private
pools. Dynamic shapes, data-dependent control flow, allocations, and
unsupported operations can limit capture. Treat ``compiled'' and
``graphed'' as measured variants, not upgrades by definition.

\subsubsection{5.7 Level 6: write a custom kernel only for a proven
target}\label{57-level-6-write-a-custom-kernel-only-for-a-proven-target}

A custom CUDA, Triton, or CUTLASS kernel is justified when profiling
identifies a dominant operation with a stable contract that existing
primitives execute poorly. Good targets often combine several of:

\begin{itemize}
\tightlist
\item
  repeated reads/writes of short-lived intermediates;
\item
  host-controlled loops over many small rows;
\item
  irregular but bounded reductions;
\item
  an operation whose semantic structure permits a single owning block;
  or
\item
  a major matrix shape that generic dispatch consistently mishandles.
\end{itemize}

The kernel needs a correctness oracle, adversarial boundary tests, dtype
and architecture coverage, resource/occupancy analysis, and an
end-to-end retention threshold. A microbenchmark win is necessary
evidence, not sufficient evidence.

\subsection{6. Batching, vectorization, fusion, and shape
design}\label{6-batching-vectorization-fusion-and-shape-design}

These words are often used interchangeably even though they solve
different problems.

\subsubsection{6.1 Batching shared-weight
work}\label{61-batching-shared-weight-work}

Given one weight matrix \(W\in\mathbb{R}^{D\times E}\) and a batch
\(X\in\mathbb{R}^{B\times D}\), ordinary batching computes

\[
Y=XW.
\]

Increasing \(B\) produces a larger GEMM and amortizes model-call
overhead. The examples share weights. This is the normal data-batch
axis.

\subsubsection{6.2 Batching independent
replicas}\label{62-batching-independent-replicas}

For \(K\) independent models,

\[
X\in\mathbb{R}^{K\times B\times D},\qquad
W\in\mathbb{R}^{K\times D\times E},
\]

and

\[
Y_k=X_kW_k.
\]

This is a batch of independent GEMMs. Flattening \(K\) into the data
batch and using one shared \(W\) would be mathematically wrong. Valid
implementations include batched matrix multiplication, grouped GEMM, a
functional vectorization such as \texttt{vmap}, or concurrent
independent stream work. Parameters, gradients, optimizer moments,
clipping norms, RNG, and failure state remain replica-local.

Replica batching is attractive for tiny models because it exposes more
work without increasing one experiment\textquotesingle s statistical
batch. It is also expensive: persistent model and optimizer state scales
with \(K\), while activations scale with the active cohort \(K B\). The
right \(K\) is a measured throughput/memory plateau, not ``as many as
fit.''

\subsubsection{6.3 Padding, bucketing, and grouped
operations}\label{63-padding-bucketing-and-grouped-operations}

Variable shapes force a choice:

\begin{itemize}
\tightlist
\item
  \textbf{Padding:} regular kernels, wasted arithmetic and memory.
\item
  \textbf{Bucketing:} several regular shape classes, grouping overhead
  and tails.
\item
  \textbf{Grouped/ragged kernels:} less padding, more complex dispatch
  and possibly lower kernel efficiency.
\end{itemize}

Measure both physical and useful work:

\[
\text{padding efficiency}
=\frac{\text{valid elements or tokens}}{\text{physically processed elements or tokens}}.
\]

A padded kernel can win even with lower padding efficiency because its
GEMMs are much better. A grouped kernel earns its complexity only if
tail waste is a material fraction of end-to-end time.

\subsubsection{6.4 Operator fusion}\label{64-operator-fusion}

Suppose three pointwise operations each read and write an \(N\)-element
tensor. Separate kernels may make several DRAM round trips and incur
three launches. A fused kernel can load once, compute in registers, and
store once. That reduces both memory traffic and launch count.

Fusion is less obviously helpful when it:

\begin{itemize}
\tightlist
\item
  creates register pressure and spills;
\item
  combines stages with incompatible parallel structure;
\item
  prevents a highly optimized library call;
\item
  lengthens live ranges and increases memory; or
\item
  explodes the number of compiled shape variants.
\end{itemize}

Fuse data movement and cheap elementwise work around dominant operators
first. Do not attempt to replace cuBLAS or a fused attention
implementation without strong shape-specific evidence.

\subsubsection{6.5 Model shape is a systems
parameter}\label{65-model-shape-is-a-systems-parameter}

Two models with equal parameter count can have very different
efficiency. Tensor Cores prefer certain dimension multiples and
sufficiently large matrix tiles. Many narrow projections may launch more
kernels and achieve less reuse than fewer well-shaped ones. Attention
cost grows approximately quadratically with sequence length for dense
attention:

\[
\text{score work}=O(BHN^2D_h),
\]

while projection work is closer to \(O(BND^2)\). A minority of long
sequences can dominate runtime.

Changing the model solely for speed is a scientific change. It should be
factorized from an execution comparison: first compare execution
strategies on the same model, then evaluate a smaller or differently
shaped architecture with its own quality gate.

\subsection{7. Numerical formats and mixed
precision}\label{7-numerical-formats-and-mixed-precision}

A floating-point format trades exponent range, significand precision,
storage, and arithmetic throughput.

{\def\LTcaptype{none} % do not increment counter
\begin{longtable}[]{@{}lrl@{}}
\toprule\noalign{}
Format & Typical storage & Main intuition \\
\midrule\noalign{}
\endhead
\bottomrule\noalign{}
\endlastfoot
FP64 & 8 bytes & large precision/range; expensive on many GPUs \\
FP32 & 4 bytes & conventional training/control precision \\
TF32 & internal Tensor Core mode & FP32 range with reduced multiply
precision \\
FP16 & 2 bytes & more significand bits than BF16, much smaller exponent
range \\
BF16 & 2 bytes & FP32-like exponent range, fewer significand bits \\
FP8 & 1 byte & very high throughput, requires careful scaling and
support \\
\end{longtable}
}

Exact bit layouts and behavior are more nuanced, but one practical
contrast is central: FP16 more readily overflows or underflows because
of its exponent range, while BF16 preserves roughly
FP32\textquotesingle s dynamic range at lower mantissa precision. This
often makes BF16 operationally convenient for training on hardware that
supports it.

\subsubsection{7.1 Mixed precision is a policy, not a dtype
switch}\label{71-mixed-precision-is-a-policy-not-a-dtype-switch}

A robust training policy may use:

\begin{itemize}
\tightlist
\item
  BF16/FP16 inputs and activations for Tensor Core operations;
\item
  FP32 accumulation where the library supports it;
\item
  FP32 loss reductions, norms, or numerically sensitive probabilities;
\item
  FP32 master parameters and optimizer moments; and
\item
  explicit conversion only at declared boundaries.
\end{itemize}

Autocast chooses dtypes for supported operations; it does not decide the
scientific tolerances or storage ownership of the application. A tensor
written into a preallocated BF16 arena must be intentionally converted
at that semantic boundary. Letting an arbitrary dtype mismatch cast
silently can hide a bug in an integer address, mask, or ownership field.

\subsubsection{7.2 Numerical equivalence cannot mean bit
identity}\label{72-numerical-equivalence-cannot-mean-bit-identity}

Floating-point addition is not associative:

\[
(a+b)+c \ne a+(b+c)
\]

in general. Parallel reductions, fused instructions, Tensor Cores, and
compiler reordering can alter roundoff. Correctness gates should
combine:

\begin{itemize}
\tightlist
\item
  exact comparisons for integer state, addresses, masks, samples, and
  RNG;
\item
  tolerance-based comparisons for outputs and losses;
\item
  relative \(L_2\) error and cosine similarity for large gradients;
\item
  finite-value and range invariants;
\item
  one-step optimizer-state comparison; and
\item
  longer replay drift bounds.
\end{itemize}

Thresholds must be fixed before selecting a faster candidate. Loosening
them after observing a failure turns validation into rationalization.

\subsubsection{7.3 Hardware support is
non-negotiable}\label{73-hardware-support-is-non-negotiable}

Turing Tensor Cores support FP16 matrix inputs but not BF16. Native BF16
Tensor Core support arrived with Ampere and is present on Hopper. A BF16
production path can therefore be tested semantically only on supporting
hardware. An RTX 2080 Ti FP32 smoke can validate control flow, storage,
CUDA residency, and many invariants, but it cannot validate H100 BF16
dispatch or numerical behavior.

\subsubsection{7.4 A worked mixed-precision training
step}\label{74-a-worked-mixed-precision-training-step}

Consider a linear layer with a million input activations. Storing those
activations in FP32 requires about 4 MB; BF16 requires about 2 MB.
Saving the input for backward therefore halves that particular
activation buffer, and a Tensor Core matrix multiplication may become
much faster. It does \textbf{not} follow that the complete training step
uses half the memory. FP32 master weights, FP32 gradients,
Adam\textquotesingle s two FP32 moment vectors, allocator workspaces,
and non-floating metadata may be unchanged.

A defensible step might look conceptually like this:

\begin{verbatim}
BF16 input and saved activations
  -> BF16 matrix operands with FP32 accumulation
  -> FP32 logits for a sensitive normalization
  -> FP32 scalar loss and gradient-norm reduction
  -> BF16/FP32 backward selected per operator
  -> FP32 gradient clipping and Adam state transition
\end{verbatim}

The cast locations are part of the interface. For example, suppose a
producer builds an FP32 feature row and the next owner is a declared
BF16 model-input arena. Converting exactly at that publication boundary
is meaningful: the address-like integer fields remain exact, and
downstream code can rely on the arena\textquotesingle s dtype.
Converting the entire record indiscriminately would corrupt its type
contract; refusing the declared floating conversion would make the mixed
precision design unusable.

FP16 and BF16 also illustrate why ``16 bit'' is not a complete
specification. The largest finite FP16 value is 65,504, whereas BF16 has
approximately the same exponent range as FP32. BF16 consequently
tolerates a much wider scale range, but its coarser significand means
nearby large values can round to the same number. The right question is
not which format is universally accurate; it is which operations need
range, which need fine relative precision, and which error the objective
and optimizer can tolerate.

\subsection{8. Asynchrony: streams, events, graphs, and accidental
barriers}\label{8-asynchrony-streams-events-graphs-and-accidental-barriers}

CUDA kernel launches are normally asynchronous with respect to the host:
the CPU enqueues work and continues. A \textbf{stream} is an ordered
sequence of device operations. Operations within a stream execute in
order; operations in different streams may overlap when dependencies and
resources permit.

\subsubsection{8.1 Events express device
dependencies}\label{81-events-express-device-dependencies}

A CUDA event can be recorded after one operation and waited on by
another stream:

\begin{verbatim}
model stream:  [forward] --record ready-- [next forward]
commit stream:               wait ready-- [publish]
\end{verbatim}

This allows a precise dependency without stopping the whole device or
forcing the CPU to wait. Device-wide synchronization is a much larger
hammer. It is appropriate at measurement, report, checkpoint, teardown,
or explicit failure boundaries; it is usually suspicious inside a
fine-grained hot loop.

\subsubsection{8.2 Concurrency is not guaranteed by using many
streams}\label{82-concurrency-is-not-guaranteed-by-using-many-streams}

Two kernels overlap only if:

\begin{itemize}
\tightlist
\item
  dependencies allow it;
\item
  they are actually enqueued early enough;
\item
  together they fit available SM/register/shared-memory resources;
\item
  no implicit synchronization intervenes; and
\item
  the device supports the needed concurrency.
\end{itemize}

A large kernel may already occupy the machine. Adding streams then
changes nothing. Conversely, concurrent tiny independent kernels can
fill otherwise idle resources. A timeline, not source-code stream count,
establishes overlap.

\subsubsection{8.3 CUDA Graphs}\label{83-cuda-graphs}

A CUDA Graph records a repeated operation/dependency schedule and
instantiates it for low-overhead replay. It is well suited to stable
shapes, addresses, and control flow. It can greatly reduce CPU launch
work for small repeated steps.

Costs and constraints include capture warmup, static buffers,
restrictions on some operations, graph-private memory pools, RNG
handling, and a cache problem when many shapes require different graphs.
A good design has a small explicit set of graph classes and an eager
fallback. Graph replay should be compared against a warmed
eager/compiled baseline with identical useful work.

\subsubsection{8.4 The hidden synchronization
checklist}\label{84-the-hidden-synchronization-checklist}

In a GPU framework, these operations often deserve inspection:

\begin{itemize}
\tightlist
\item
  converting a device scalar to a Python number;
\item
  copying a tensor to CPU;
\item
  host-side branching on device data;
\item
  exception or assertion paths that inspect device state;
\item
  allocator phase changes and cache-emptying calls;
\item
  distributed collectives in inconsistent order;
\item
  default-stream work interacting with non-default streams; and
\item
  timing without an explicit synchronization boundary.
\end{itemize}

The answer is not to ban them. It is to make them rare, intentional, and
visible in a trace.

\subsubsection{8.5 A worked two-stream
pipeline}\label{85-a-worked-two-stream-pipeline}

Suppose batch \(i\) needs 3 ms of model work and 1 ms of publication
work. A single stream executes

\begin{verbatim}
model(i) [---] publish(i) [-] model(i+1) [---] publish(i+1) [-]
\end{verbatim}

for a nominal 8 ms across two batches. If publication uses few enough
resources and batch \(i+1\) is independent, two streams can instead
express

\begin{verbatim}
model stream:   model(i) [---] model(i+1) [---]
commit stream:               publish(i) [-]   publish(i+1) [-]
\end{verbatim}

An event recorded after \texttt{model(i)} permits \texttt{publish(i)} to
start without a device-wide barrier. The first publication may overlap
the second model, so the idealized makespan falls toward 7 ms. It cannot
fall to 6 ms merely because two streams exist: each stream preserves its
own order, publication depends on its model, and the two kernels may
contend for the same resources.

Now add \texttt{loss\ =\ device\_loss.item()} between the two model
launches. The host must wait until the first loss is available before it
can enqueue later work; the intended overlap may vanish. This is why
source-level asynchrony is not enough. The proof is a timeline showing
early enqueue, event dependencies, and actual overlap on the target
device.

\subsection{9. Memory engineering as system
design}\label{9-memory-engineering-as-system-design}

GPU memory is not just a bag into which tensors are placed. It is an
ownership and scheduling constraint. A reliable accelerator system
answers four questions for every material allocation:

\begin{enumerate}
\def\labelenumi{\arabic{enumi}.}
\tightlist
\item
  Who owns it?
\item
  What exact shape and dtype may it have?
\item
  What is its last consumer?
\item
  What happens if the producer fails or exceeds capacity?
\end{enumerate}

\subsubsection{9.1 Static arenas versus dynamic tensor
allocation}\label{91-static-arenas-versus-dynamic-tensor-allocation}

Dynamic framework allocation is convenient and often entirely adequate.
For a large irregular pipeline, however, thousands of differently shaped
allocations make capacity hard to prove and failure hard to make atomic.
A fixed arena can instead provide typed views at aligned offsets and
monotone row reservations:

\begin{verbatim}
reserve [begin, end) -> fill every declared field -> publish tail
\end{verbatim}

Readers observe only the published prefix. If filling fails, the
unpublished reservation can be discarded. This resembles a database
transaction more than a conventional Python list.

Advantages include predictable addresses, exact analytical storage,
fewer allocator calls, easy high-water accounting, and graph-friendly
stability. Costs include internal slack, up-front capacity work, more
rigid schemas, and the need for explicit lifetime discipline. Static
storage is high ROI only when the workload has useful bounds and
correctness benefits from the contract.

\subsubsection{9.2 Capacity and microbatching are different
dimensions}\label{92-capacity-and-microbatching-are-different-dimensions}

This distinction prevents a common failure. A \textbf{cumulative
capacity} covers everything retained over a phase: for example, one row
per root across an entire rollout. A \textbf{workspace or pack capacity}
covers only the rows processed concurrently: for example, 256 roots in
one model call.

If cumulative storage is accidentally sized to one pack, several packs
may run successfully before deterministic overflow. If workspace is
sized to the full cumulative count, the program may waste tens of
gigabytes. Each bound should be named by lifetime:

\[
\text{cumulative rows} \ne \text{maximum simultaneous rows}.
\]

\subsubsection{9.3 Preflight admission}\label{93-preflight-admission}

A capacity compiler can transform a small workload description into all
derived row counts, tensor shapes, workspaces, and phase peaks. It
should reject an impossible run before expensive rollout begins.

For each phase \(\phi\), a conservative admission inequality is

\[
M_{\mathrm{persistent}}
+M_{\mathrm{live},\phi}
+M_{\mathrm{workspace},\phi}
+M_{\mathrm{allocator/graph},\phi}
+M_{\mathrm{reserve}}
\leq M_{\mathrm{device}}.
\]

The reserve is not a substitute for missing formulas. It covers residual
uncertainty, library workspaces, fragmentation, and operational safety.
Compare the analytical prediction with measured peaks and refine the
model using held-out configurations. Do not fit a line to one point and
call it a proof.

\subsubsection{9.4 Failure atomicity}\label{94-failure-atomicity}

Accelerator failure can arrive late: a nonfinite loss after forward, an
OOM in backward, a distributed peer failure, or a capacity overflow
after partial work. If some replicas update parameters and others do
not, an independent-experiment batch is corrupted. If queue tails
publish before payloads are complete, consumers see invalid state.

Transactional design separates:

\begin{itemize}
\tightlist
\item
  provisional work;
\item
  validation/failure agreement;
\item
  committed optimizer or storage state; and
\item
  durable run/checkpoint publication.
\end{itemize}

Rollback state costs memory, so it belongs in the capacity model. A
``preflight'' that recomputes the full learner merely to avoid rollback
may be correct yet double the dominant work; the system should measure
that cost explicitly.

\subsection{10. The software stack: where CUDA performance decisions are
made}\label{10-the-software-stack-where-cuda-performance-decisions-are-made}

Application code normally sits above several layers:

\begin{verbatim}
Python model / simulation
  -> PyTorch dispatcher and autograd
    -> compiler/vectorization layers
      -> ATen operators
        -> cuBLAS, cuDNN, custom CUDA/Triton/CUTLASS kernels
          -> CUDA runtime/driver
            -> GPU
\end{verbatim}

A single line of Python may launch one highly optimized kernel, dozens
of small kernels, or a fallback implementation. Conversely, a Python
loop can be cheap if each iteration launches a large saturating
operation. Source-line count is not a performance model.

\subsubsection{10.1 Eager execution}\label{101-eager-execution}

PyTorch eager mode dispatches operations as Python executes them and
constructs an autograd graph for differentiable work. It is productive,
observable, and a valuable correctness oracle. Its weaknesses appear
when many small operations create launch gaps or write materialized
intermediates that a compiler could fuse.

\subsubsection{10.2 Autograd and activation
ownership}\label{102-autograd-and-activation-ownership}

Backward needs saved forward intermediates. The apparent forward working
set therefore understates training memory. Activation checkpointing
trades extra forward recomputation for lower retained activation memory.
It is useful when memory is binding and recomputation is cheaper than a
smaller useful batch, but harmful if launch or compute time is already
the main limit.

An accidental \texttt{no\_grad}, in-place mutation of a saved tensor,
detached copy, or incorrect functional parameter view is a correctness
error. Vectorized replica training needs gradient tests at the flat
parameter level, not only matching forward outputs.

\subsubsection{\texorpdfstring{10.3 Functional transforms and
\texttt{vmap}}{10.3 Functional transforms and vmap}}\label{103-functional-transforms-and-vmap}

A module normally owns one parameter set. Functionalization makes
parameters explicit inputs. Stacking \(K\) parameter sets and applying
\texttt{vmap} can ask the framework to vectorize a function over the
replica axis while retaining one canonical model definition.

This is elegant when every operator has an efficient batching rule. A
fallback may instead loop or introduce awkward layouts. Attention
backward, scatter, indexing, mutation, and per-replica randomness
deserve direct inspection. A successful \texttt{vmap} call establishes
semantics, not efficient dispatch.

\subsubsection{10.4 Explicit batched and grouped linear
algebra}\label{104-explicit-batched-and-grouped-linear-algebra}

When vectorization dispatch is poor, an explicit replica-native path can
store linear weights as \texttt{{[}K,\ out,\ in{]}} and use BMM/grouped
GEMM. This gives direct control over strides, padding, and the folding
of replica and example axes into attention batches. It also duplicates
implementation surface and complicates checkpoint conversion.

The clean compromise is one scientific model/loss contract with a
narrow, tested execution backend. Keep the ordinary module as a
correctness oracle; retain the specialized backend only if it produces a
material full-step gain.

\subsubsection{\texorpdfstring{10.5
\texttt{torch.compile}}{10.5 torch.compile}}\label{105-torchcompile}

\texttt{torch.compile} captures compatible PyTorch regions and uses a
backend such as Inductor to generate optimized code. Potential benefits
include pointwise fusion, lower dispatch overhead, layout changes, and
CUDA Graph integration.

Watch for:

\begin{itemize}
\tightlist
\item
  graph breaks;
\item
  compilation repeated for new shapes or scalar values;
\item
  cold-start time included in steady-state benchmarks;
\item
  numerical changes from reordering;
\item
  larger code/workspace caches; and
\item
  fallback operators dominating the supposedly compiled region.
\end{itemize}

Compile a bounded learner replay before an entire dynamic simulation.
Stable leaf functions are much easier to reason about than
data-dependent queue orchestration.

\subsubsection{10.6 Custom extensions}\label{106-custom-extensions}

A C++/CUDA extension can combine a host loop into one compiled launch
sequence or implement an operation absent from PyTorch. Its build is
architecture-specific state. Record toolkit, compiler, ABI, PyTorch
version, target compute capabilities, source revision, and cache
directory. A cache key that omits any of these can reuse an incompatible
binary.

\subsection{11. Measurement and profiling: how to learn what the machine
is
doing}\label{11-measurement-and-profiling-how-to-learn-what-the-machine-is-doing}

There are three complementary levels of evidence.

\subsubsection{11.1 Application metrics}\label{111-application-metrics}

Application metrics answer whether the system accomplished the intended
work:

\begin{itemize}
\tightlist
\item
  useful simulations or decisions per second;
\item
  useful learner presentations per second;
\item
  optimizer transitions per second;
\item
  time-to-quality or time-to-statistical precision;
\item
  exact sample count and padding fraction; and
\item
  phase time and memory.
\end{itemize}

These decide whether an optimization matters. Hardware counters explain
why.

\subsubsection{11.2 Timeline profiling}\label{112-timeline-profiling}

Nsight Systems or a framework trace shows CPU threads, CUDA API calls,
kernels, copies, streams, events, and communication over time. It
answers:

\begin{itemize}
\tightlist
\item
  Is the GPU idle between kernels?
\item
  Is the host submitting work quickly enough?
\item
  Are supposedly concurrent lanes actually overlapping?
\item
  Where are device-wide synchronizations?
\item
  Do H2D/D2H copies sit on the critical path?
\item
  Are distributed ranks waiting on collectives?
\item
  Which phase owns the long interval?
\end{itemize}

Annotate domain phases with NVTX ranges. A raw sequence of kernel names
is much harder to map back to the system.

\subsubsection{11.3 Kernel profiling}\label{113-kernel-profiling}

Nsight Compute examines selected kernels. Useful questions include:

\begin{itemize}
\tightlist
\item
  achieved versus shape-specific compute throughput;
\item
  HBM/L2/L1 traffic and hit rates;
\item
  occupancy and the resource limiting resident blocks;
\item
  warp execution efficiency and divergence;
\item
  memory-load coalescing;
\item
  instruction mix and issue stalls; and
\item
  whether Tensor Core instructions were actually issued.
\end{itemize}

Profilers perturb execution and can serialize kernels. Use separate
unprofiled timing runs and bounded profiling captures. Do not feed
profiler-perturbed wall times into the main throughput result.

\subsubsection{\texorpdfstring{11.4 \texttt{nvidia-smi} is operational
telemetry}{11.4 nvidia-smi is operational telemetry}}\label{114-nvidia-smi-is-operational-telemetry}

\texttt{nvidia-smi} is excellent for confirming the selected GPU, memory
use, power, temperature, clocks, compute processes, ECC state, and
whether another job is present. Its utilization percentage is a coarse
sampling signal. It cannot prove Tensor Core efficiency, identify a
memory bottleneck, or validate Flash attention dispatch.

\subsubsection{11.5 A trustworthy benchmark
protocol}\label{115-a-trustworthy-benchmark-protocol}

A performance comparison should:

\begin{enumerate}
\def\labelenumi{\arabic{enumi}.}
\tightlist
\item
  pin exact hardware, software, source, config, and data/replay
  identity;
\item
  run on an idle exclusive device when possible;
\item
  separate initialization/compile/autotune from steady state;
\item
  synchronize immediately around the timed boundary;
\item
  compare identical useful work;
\item
  alternate candidates to reduce drift and thermal/order bias;
\item
  collect enough repeats to show variability;
\item
  report median and dispersion, not only the best run;
\item
  retain raw machine-readable artifacts; and
\item
  verify clean device teardown.
\end{enumerate}

For expensive rentals, prepare a resumable experiment matrix and every
command before the rental begins. Write progress atomically after each
completed unit. A timeout with partial evidence is much better than an
empty log.

\subsection{12. Correctness, determinism, and
portability}\label{12-correctness-determinism-and-portability}

Performance engineering changes schedules. Schedules influence RNG
consumption, floating-point reduction order, allocator state, and
failure timing. Correctness cannot be bolted on after the fast path
exists.

\subsubsection{12.1 Logical randomness}\label{121-logical-randomness}

If a random draw depends on which queue happens to run first, changing a
microbatch or stream schedule changes the experiment. A stronger design
derives randomness from a logical address such as

\[
(\text{seed},\text{iteration},\text{root},\text{node},\text{purpose},\text{attempt}).
\]

Then physical batching may change without changing logical samples.
Explicit per-replica generators are another workable design when their
consumption order is part of the contract. The key is to decide,
document, and test the invariance boundary.

\subsubsection{12.2 Determinism has
levels}\label{122-determinism-has-levels}

Useful reproducibility profiles include:

\begin{itemize}
\tightlist
\item
  \textbf{bitwise:} same bits on a fixed software/hardware stack;
\item
  \textbf{seeded/numerical:} same logical samples with bounded floating
  drift;
\item
  \textbf{throughput:} allows nondeterministic fast kernels while
  preserving statistical and state contracts.
\end{itemize}

Demanding bit identity everywhere can reject valuable optimized
reductions. Demanding no identity anywhere makes debugging scientific
regressions nearly impossible. Different workflows can select different
profiles explicitly.

\subsubsection{12.3 Architecture
portability}\label{123-architecture-portability}

A test on an older GPU can validate shared CUDA logic while missing:

\begin{itemize}
\tightlist
\item
  unsupported numeric formats;
\item
  different library dispatch;
\item
  architecture-specific instructions;
\item
  larger-memory geometry;
\item
  race timing exposed only at higher concurrency;
\item
  new framework compatibility defects; and
\item
  native-extension build errors for the new SM target.
\end{itemize}

Portability therefore uses a ladder: fast local tests on the available
device, architecture-specific smoke tests, then production-scale tests
on the target. ``Works on CUDA'' is not a single binary property.

\subsection{13. Multi-GPU systems in one
page}\label{13-multi-gpu-systems-in-one-page}

Multiple GPUs introduce a network and a distributed state machine.
CUDA-aware collectives are commonly provided by NCCL:

\begin{itemize}
\tightlist
\item
  all-reduce combines gradients or statistics across ranks;
\item
  broadcast distributes one rank\textquotesingle s data;
\item
  all-gather collects shards from every rank;
\item
  reduce-scatter combines and partitions data; and
\item
  point-to-point operations express explicit exchange.
\end{itemize}

For data-parallel training, a rough step model is

\[
T_{\rm step}
\approx
T_{\rm local\ compute}
+T_{\rm exposed\ communication}
+T_{\rm imbalance}
+T_{\rm synchronization}.
\]

Communication may overlap backward compute, but small models often have
too little compute to hide collective latency. All ranks must enter
collectives in compatible order. One rank failing while others wait can
look like a hang.

Scale only after one GPU is efficient enough and the scientific
global-batch semantics are explicit. Compare strong scaling (same total
work, more GPUs) and weak scaling (more total work with GPU count)
separately. Aggregate throughput can rise while per-item latency and
cost efficiency worsen.

\subsection{14. A worked diagnosis: why a tiny transformer disappoints
on a huge
GPU}\label{14-a-worked-diagnosis-why-a-tiny-transformer-disappoints-on-a-huge-gpu}

Imagine a transformer whose one training step contains many small
projections, normalizations, gathers, and optimizer operations. Moving
it from an RTX 2080 Ti to an H100 gives far more memory bandwidth and
Tensor Core peak, yet observed throughput improves only modestly.

Start with hypotheses, not code:

\begin{enumerate}
\def\labelenumi{\arabic{enumi}.}
\tightlist
\item
  \textbf{Not enough work:} GEMMs are too small to fill H100 Tensor
  Cores.
\item
  \textbf{Launch bound:} the CPU submits hundreds of short kernels with
  gaps.
\item
  \textbf{Memory bound:} gathers and pointwise operators dominate matrix
  work.
\item
  \textbf{Bad dispatch:} attention uses a math fallback or incompatible
  layout.
\item
  \textbf{Serialization:} independent replicas run one after another.
\item
  \textbf{Recomputation:} a safety preflight performs forward/backward
  twice.
\item
  \textbf{Measurement error:} cold compilation or different useful work
  pollutes the comparison.
\end{enumerate}

The experiment sequence should be:

\begin{enumerate}
\def\labelenumi{\arabic{enumi}.}
\tightlist
\item
  Reconcile exact work and isolate learner time.
\item
  Capture a timeline for launch gaps and serialization.
\item
  Record shapes and prove attention/GEMM dispatch.
\item
  Increase the cheapest honest parallel axis until throughput plateaus.
\item
  Compile or graph a fixed replay if launch gaps remain.
\item
  Vectorize optimizer control if it remains visible.
\item
  Try replica-axis \texttt{vmap}, explicit BMM, or stream waves against
  the same serial replay.
\item
  Profile only dominant remaining kernels.
\item
  Write a custom kernel only if its Amdahl ceiling is worthwhile.
\end{enumerate}

This ordering distinguishes ``the GPU needs more work'' from ``each unit
of work is inefficient.'' Both can be true, but they require different
fixes.

\subsection{15. A compact Pareto
checklist}\label{15-a-compact-pareto-checklist}

Before optimizing:

\begin{itemize}
\tightlist
\item
  What is the useful-work numerator?
\item
  What must remain exact, and what may differ within tolerance?
\item
  What percentage of end-to-end time can the proposal affect?
\item
  Is the target launch-, memory-, compute-, communication-, or
  capacity-bound?
\end{itemize}

Before increasing batch or concurrency:

\begin{itemize}
\tightlist
\item
  Is the axis scientifically independent or correctly normalized?
\item
  Which state scales with \(K\), \(B\), sequence length, or sample
  count?
\item
  Are cumulative and simultaneous capacities separate?
\item
  Does throughput plateau before memory fills?
\end{itemize}

Before using mixed precision:

\begin{itemize}
\tightlist
\item
  Does the exact GPU support it?
\item
  Which state remains FP32?
\item
  Where are casts allowed?
\item
  What are the predeclared numerical gates?
\end{itemize}

Before compiling or graphing:

\begin{itemize}
\tightlist
\item
  Are shapes and addresses stable?
\item
  Are graph breaks/recompilations bounded?
\item
  Is cold-start time reported separately?
\item
  Is private-pool/cache memory admitted?
\end{itemize}

Before writing a custom kernel:

\begin{itemize}
\tightlist
\item
  Is the operation dominant in a real trace?
\item
  What is the maximum end-to-end speedup?
\item
  Is there a simple trusted oracle?
\item
  Can a library, fusion, layout, or batching change solve it first?
\item
  What measured retention threshold will justify maintenance?
\end{itemize}

\begin{center}\rule{0.5\linewidth}{0.5pt}\end{center}

\section{Part II --- How this project uses those
ideas}\label{part-ii--how-this-project-uses-those-ideas}

\subsection{16. The computational problem in this
repository}\label{16-the-computational-problem-in-this-repository}

This repository trains strategies for configurable no-limit Texas
Hold\textquotesingle em cash poker, an imperfect-information game. The
active game has three per-run dials:

\begin{itemize}
\tightlist
\item
  players \(P\in\{2,\ldots,6\}\);
\item
  equal stack depth \(S\in[8,128]\) big blinds in 0.5 bb increments; and
\item
  retained ranks \(R\in\{4,\ldots,13\}\) with all four suits.
\end{itemize}

It requires \(2P+5<4R\), uses 0.5/1 small/big blinds, no ante or rake,
and equal stacks. The documented maximal player-node count in one
trajectory is

\[
4P+(S-1)(P-1)=SP+3P-S+1.
\]

Unlike training on a static image dataset, one iteration constructs
stochastic game trajectories, performs policy inference at decision
nodes, handles chance and terminal states, evaluates hidden-card
outcomes, propagates values through a topology, creates aligned training
targets, and only then performs learner updates. Shapes vary by street,
public history, legal action count, hidden-hand domain, rollout mode,
and estimator.

The central active pipeline is:

\begin{verbatim}
strict config and immutable capacity plan
  -> fixed accelerator-resident storage owners
  -> family-scheduled rollout
       model -> sample -> transition -> atomic publication
  -> sampled showdown or Rao--Blackwell completion
  -> compact topology and reverse layers
  -> layered target propagation
  -> aligned learner sampling
  -> forward/loss/backward and dense FP32 optimizer transition
  -> metrics and atomic checkpoint
\end{verbatim}

CUDA supports every production workflow. MPS is limited to supported
inference paths. CPU work is intentionally small orchestration such as
config parsing, hashing, checkpoint I/O, topology metadata, and report
formatting; it does not own production rollout tensors, model forwards,
loss/backward, posterior work, or showdown evaluation.

\subsubsection{16.1 Poker and information sets in one
page}\label{161-poker-and-information-sets-in-one-page}

In heads-up no-limit Hold\textquotesingle em, each player receives two
private cards. Up to five public board cards arrive over the flop, turn,
and river. On each betting street a player may fold, check/call when
legal, or raise to a legal amount. The hand ends by a fold or by
comparing five-card poker hands at showdown. Stacks, bets, board cards,
and action history are public; hole cards are not.

That hidden information changes the learning object. A policy cannot
condition on the opponent\textquotesingle s actual private cards,
because the player does not observe them. It conditions on an
\textbf{information set}: everything the acting player can know,
together with a belief or reach distribution over compatible hidden
hands. Two world states that differ only in an unseen opponent hand
belong to the same public decision problem. Revealing that hand to the
network during training would create a policy that cannot be executed at
play time.

The project consequently carries vectors over private-hand domains.
Before the flop, a full 52-card deck has 169 strategically distinct
rank/suitedness classes: 13 pairs, 78 suited non-pairs, and 78 offsuit
non-pairs. After public cards appear, suit-specific collisions matter,
so the natural domain is the \(\binom{52}{2}=1,326\) concrete two-card
combinations. The network produces policy/value information over these
domains and the learner selects or weights rows according to the
estimator contract.

This explains why one ``poker example'' is much larger than one scalar
feature row. It contains public geometry, action history, and
distributions over many possible private states. It also explains the
strict evaluator boundary: showdown and completion are large
hidden-state calculations, not incidental Python rules.

\subsubsection{16.2 A concrete HU 32 bb shape
calculation}\label{162-a-concrete-hu-32-bb-shape-calculation}

Take the selected baseline \(P=2\), \(S=32\), \(R=13\), with two learned
raise components. The repository\textquotesingle s geometry formulas
give:

{\def\LTcaptype{none} % do not increment counter
\begin{longtable}[]{@{}lrr@{}}
\toprule\noalign{}
Quantity & Calculation & Result \\
\midrule\noalign{}
\endhead
\bottomrule\noalign{}
\endlastfoot
Deck cards & \(4R\) & 52 \\
Concrete private hands & \(\binom{52}{2}\) & 1,326 \\
Preflop private classes & \(R^2\) & 169 \\
Maximum player nodes & \(SP+3P-S+1\) & 39 \\
History tokens & 39 actions + 3 street separators & 42 \\
Public feature width & \(8+10P\) & 28 \\
History feature width & \(10\cdot42\) & 420 \\
Retained preflop row width & \(28+2\cdot169+420\) & 786 \\
Retained postflop row width & \(28+2\cdot1,326+420\) & 3,100 \\
\end{longtable}
}

The maximum structural trajectory has 39 player nodes, three
chance/street nodes, and one leaf: 43 total nodes before topology
compaction. These are capacity bounds, not predictions of the average
hand length.

The transformer sequence length is different from the retained flat
feature width. Under the pilot configuration, the conservative preflop
maximum is

\[
2(169)+150+1+42+5=536\ \text{tokens},
\]

where 150 accounts for seat, scalar, event, and half-blind-tick
features; the remaining terms are a global token, public history, and
five configured action tokens. The corresponding current full-row
postflop maximum is

\[
2(1,326)+150+6+42+5=2,855\ \text{tokens}.
\]

On a flop, only \(\binom{49}{2}=1,176\) concrete hands avoid the three
board cards. Thus 150 of the 1,326 rows are logically impossible for
each example. The model has compatibility maps and compact
gather/scatter machinery, but the conservative queued call currently
retains the full padded width unless a fixed compact size is supplied.
This is an excellent candidate for a measured optimization: it offers a
clear \(1,326/1,176\approx1.13\) row ratio, but only a trace can show
whether removing those rows materially reduces end-to-end time.

These calculations illustrate three layers of ``shape'' that must not be
confused:

\begin{itemize}
\tightlist
\item
  cumulative trajectory/topology capacity;
\item
  retained storage width for one public state; and
\item
  physical model sequence length after tokenization and padding.
\end{itemize}

Changing any one may leave the others unchanged.

\subsubsection{16.3 The mode space}\label{163-the-mode-space}

The core generates twelve mode contracts from

\begin{verbatim}
{self-play, best response}
× {sequential, branched fixed, branched KKT}
× {sampled, Rao--Blackwell}
\end{verbatim}

Normal sequential training samples a concrete private outcome.
Sequential RB training integrates additional hidden-state uncertainty
through completion and mass/moment targets. Branched modes deliberately
expand selected decisions. These variants share storage, rollout,
target, and optimizer authorities while changing declared fields,
multiplicity, estimator logic, and presentation.

\subsubsection{16.4 Sampled versus Rao-\/-Blackwell values: a toy
example}\label{164-sampled-versus-rao--blackwell-values-a-toy-example}

Suppose a hidden opponent state is \(A\) with probability 0.6 and \(B\)
with probability 0.4. For a fixed public decision, imagine the terminal
contribution is \(+1\) in state \(A\) and \(-1\) in state \(B\). A
Normal sampled target draws one state. Its expectation is

\[
0.6(1)+0.4(-1)=0.2,
\]

but one sample is always \(+1\) or \(-1\), with variance
\(1-0.2^2=0.96\). An exact conditional expectation would publish \(0.2\)
directly and remove that hidden-state sampling variance. This is the
intuition behind Rao-\/-Blackwellization. More generally,

\[
\operatorname{Var}(X)
=\mathbb{E}[\operatorname{Var}(X\mid Y)]
+\operatorname{Var}(\mathbb{E}[X\mid Y]),
\]

and replacing \(X\) by \(\mathbb{E}[X\mid Y]\) removes the first
nonnegative term.

The real project cannot enumerate every relevant completion for free.
Its RB path constructs posterior/profile samples and mass/moment
targets, then asks the evaluator-owned completion operation for the
relevant hand-rank work. The selected H100 baseline used 1,024
completion samples. RB therefore trades more workspace and evaluator
work for a potentially lower-variance target; changing that sample count
is a scientific estimator change, not merely a batch-size tweak.

\subsubsection{16.5 One decision node through the
system}\label{165-one-decision-node-through-the-system}

For a concrete mental execution, follow one postflop public node:

\begin{enumerate}
\def\labelenumi{\arabic{enumi}.}
\tightlist
\item
  The fixed state owner holds stacks, commitments, board, statuses,
  reaches, and packed history on CUDA.
\item
  The scheduler places the row in a semantic family such as
  \texttt{post\_2}, so rows with a compatible model contract can be
  executed together.
\item
  The model reconstructs public features and compatibility masks,
  evaluates attention over the private-hand representation, and emits
  value plus factorized fold/check-call/raise information for all 1,326
  output rows.
\item
  A logical RNG address selects the action required by the estimator
  without making randomness depend on queue order.
\item
  The transition writes the child state into an unpublished reservation.
  Only after every field succeeds does the transaction publish the new
  tail.
\item
  A terminal child becomes a leaf. Sampled mode evaluates its sampled
  hidden state; RB mode also constructs completion work over posterior
  uncertainty.
\item
  The CPU-owned compact topology records structural relationships, while
  numerical leaf and layer data remain accelerator-resident.
\item
  Reverse-layer CUDA propagation produces decision targets. The learner
  then samples aligned public rows, gathers the relevant private
  targets, performs forward/loss/backward, and commits one dense FP32
  Adam transition.
\end{enumerate}

Thousands of such nodes share each physical call, and phases have
different pack limits. The example is conceptually sequential, but the
implementation\textquotesingle s purpose is to batch independent rows at
every honest axis while preserving the same transaction and estimator
semantics.

\subsubsection{16.6 Why this is a difficult GPU
workload}\label{166-why-this-is-a-difficult-gpu-workload}

Several tensions are present at once:

\begin{itemize}
\tightlist
\item
  trajectories are irregular, while GPUs prefer regular batches;
\item
  public states are cheap relative to hidden-hand and completion
  domains;
\item
  preflop and postflop use materially different model/input shapes;
\item
  the simulator generates the learner\textquotesingle s data, so
  production and consumption rates must be balanced;
\item
  correctness requires schedule-independent randomness and transactional
  publication;
\item
  RB improves an estimator but introduces large sample/workspace
  dimensions;
\item
  independent replicas can expose parallelism but replicate persistent
  state; and
\item
  the neural model is small enough that a large H100 can be launch- or
  shape-limited rather than peak-compute-limited.
\end{itemize}

The project is therefore a systems pipeline with neural kernels inside
it, not merely ``a transformer on a GPU.''

\subsection{17. Where the high-ROI ideas appear in the
implementation}\label{17-where-the-high-roi-ideas-appear-in-the-implementation}

This section describes the architectural intent visible at the current
repository revision. Exact class and field names can evolve; the systems
ideas are more durable.

\subsubsection{17.1 One canonical config becomes one resolved execution
manifest}\label{171-one-canonical-config-becomes-one-resolved-execution-manifest}

The scientific config names game, mode, model, objective, optimizer,
schedule, and seed. Hardware execution config names the explicit device,
precision, attention backend, reserve, and physical microbatch limits. A
capacity compiler derives partition rows, topology and leaf bounds,
token bounds, workspaces, and phase envelopes from that canonical input
and the observed CUDA memory.

This addresses a subtle but high-value systems problem: duplicated
derived configuration is executable disagreement. If YAML says one input
width while the model constructs another, or one generator computes
folded capacity from a pack while runtime retains it for a phase, the
GPU eventually reveals the disagreement as an overflow or dtype/shape
error. Derive facts once and save the resolved manifest beside the run.

\subsubsection{17.2 Fixed SoA slabs and monotone
publication}\label{172-fixed-soa-slabs-and-monotone-publication}

\texttt{SoASchema}, \texttt{MonotoneSlab}, \texttt{ExactBlock}, and
generated owners provide fixed, aligned, accelerator-resident storage.
Typed fields are views over byte arenas. One lane reserves, fills, and
publishes monotone prefixes; capacity does not grow at runtime.

This combines performance and correctness:

\begin{itemize}
\tightlist
\item
  SoA supports coalesced field-wise kernels;
\item
  fixed addresses reduce allocation churn and suit asynchronous
  execution;
\item
  published tails are explicit producer/consumer boundaries;
\item
  high-water marks make memory use auditable;
\item
  whole failed attempts can be discarded; and
\item
  an overflow fails clearly instead of paging state through CPU memory.
\end{itemize}

The tradeoff is that geometry proofs become first-class software. The
history shows why that is worthwhile: wrong capacity is caught as a
local invariant rather than becoming silent corruption.

\subsubsection{17.3 Separate cumulative owners from bounded
workspaces}\label{173-separate-cumulative-owners-from-bounded-workspaces}

The current capacity code distinguishes persistent/cumulative storage
from bounded rollout, chance, completion, target, and learner packs. For
example, \texttt{rollout\_rows}, \texttt{completion\_rows}, and
\texttt{learner\_rows} are independent physical execution limits even
when the scientific learner batch remains 256.

This is exactly the logical/physical split from the generic section. It
permits a large statistical batch or root count while scheduling each
expensive phase with a different memory-optimal pack. It also explains
why ``batch 256'' does not prove every internal model call contains 256
rows.

\subsubsection{17.4 Explicit phase
lifetimes}\label{174-explicit-phase-lifetimes}

Rollout-only owners are released after sealing; leaves survive until
completion; layer work survives until target publication; learner heads,
targets, and indices survive through backward. Completion is
intentionally serialized after rollout in the documented design so that
evaluator workspace can reuse memory rather than compete with peak
rollout state.

This sacrifices possible overlap in exchange for a smaller and more
provable memory envelope. On an 80 GB H100, capacity often buys more
useful root-level parallelism than speculative overlap would buy
latency. The right decision is phase-specific and measured, but the
current serialization has a coherent systems rationale.

\subsubsection{17.5 Device-native rollout transactions and
streams}\label{175-device-native-rollout-transactions-and-streams}

\texttt{RolloutTransaction} is the common model/sample/transition/commit
path. The implementation uses CUDA streams and events to express
readiness and atomic publication. Logical RNG addresses separate random
meaning from queue order. Failure metadata is device-resident until an
explicit inspection boundary.

This avoids a common anti-pattern in simulators: move state to CPU after
every step, branch in Python, then upload the next batch. The CPU still
owns small scheduling and topology responsibilities, but numerical
rollout state remains accelerator-native.

\subsubsection{17.6 A compiled layered-target
kernel}\label{176-a-compiled-layered-target-kernel}

The rollout produces a compact topology. Reverse-layer target
propagation is implemented by a compiled CUDA executor that launches one
kernel per nonempty layer without a Python loop over individual
nodes/actions or an inter-layer host synchronization. One thread block
owns one parent and reduces children in logical order. Two work arenas
alternate across layers, and short branches can remain in their
producing slot until their last consumer rather than being copied.

This is a model high-ROI custom kernel target: stable semantics,
repeated bounded reductions, clear ownership, and a
host-loop/materialization cost that generic tensor expressions handled
poorly. The phase-closure audit recorded a matched old
target-propagation path at 14.670 seconds versus a 0.00576-second CUDA
interval after the redesign. That historical comparison should not be
generalized to the full application, but it illustrates the payoff of
moving the right structural loop.

\subsubsection{17.7 Optimized attention with an explicit
backend}\label{177-optimized-attention-with-an-explicit-backend}

Poker\textquotesingle s preflop and postflop models use reach-biased
multi-head attention. The runtime accepts explicit \texttt{sdpa} or
\texttt{chunked\_dense} strategies rather than an implicit ``auto''
application-level fallback. PyTorch SDPA is the high-ROI route to
optimized attention kernels; the chunked path provides a bounded
alternative when its contract is selected.

There is an important evidence boundary: an SDPA configuration proves
that the operator was requested, not that every real shape used Flash
attention. The August H100 run configured and exercised BF16 SDPA but
did not capture a formal kernel-dispatch proof. Future profiling should
close that gap before attention performance conclusions are made.

\subsubsection{17.8 The evaluator
boundary}\label{178-the-evaluator-boundary}

The sibling \texttt{poker-gpu-evaluator} contract owns showdown
evaluation and the completion operations that produce \texttt{LayerWork}
involving hand ranks. This repository owns public game state, policy,
likelihood reaches, posterior/profile sampling, and folded summaries.
There is no CPU poker-math fallback.

That is both a modularity and performance decision. One bound evaluator
per deck geometry can retain/reuse its CUDA workspace. A strict
interface prevents an innocent fallback from staging enormous
hidden-hand computation through host memory.

\subsubsection{17.9 Dense FP32 optimizer state and transactional
learning}\label{179-dense-fp32-optimizer-state-and-transactional-learning}

Model execution may use autocast BF16, while dense parameters/gradients,
optimizer controls, and moments retain FP32 semantics. The learner
transaction protects model, optimizer, RNG, replica version, and run
state so a late lane or rank failure does not partially publish an
iteration.

This costs memory. The current design chooses debuggability and exact
rollback over a smaller but fragile state machine. Optimization should
focus on batching finite checks, clipping, and Adam operations without
changing per-lane norms or reductions.

\subsubsection{17.10 A concrete replica-memory
calculation}\label{1710-a-concrete-replica-memory-calculation}

The production HU configuration contains 4,531,771 preflop and 6,602,467
postflop parameters, or 11,134,238 total. One FP32 vector over every
parameter therefore occupies

\[
4(11,134,238)=44,536,952\ \text{bytes}\approx42.47\ \text{MiB}.
\]

The current learner/reference audit expands that vector into the
following persistent ownership per independent replica:

{\def\LTcaptype{none} % do not increment counter
\begin{longtable}[]{@{}lr@{}}
\toprule\noalign{}
Persistent state & Approximate bytes \\
\midrule\noalign{}
\endhead
\bottomrule\noalign{}
\endlastfoot
Learned parameters & 44,536,952 \\
Dense gradients & 44,536,952 \\
Two Adam moments & 89,073,904 \\
Parameter and two-moment rollback snapshots & 133,610,856 \\
Frozen reference parameters & 44,536,952 \\
\textbf{Total before small buffers} & \textbf{356,295,616 (0.332
GiB)} \\
\end{longtable}
}

BF16 autocast does not halve these quantities because their declared
owner is FP32. At \(K=32\), persistent learner/reference state alone is
already about \(32(0.332)=10.6\) GiB before activations, rollout owners,
completion workspaces, targets, reference caches, allocator reserve, or
safety margin.

This is why replica count and physical learner batch must be planned
jointly. The audit\textquotesingle s empirical learner-only increment
was approximately

\[
M_{\rm active}(K,B)\approx
K\left(0.332+0.00360B\right)\ \text{GiB},
\]

where \(B\) is the simultaneously active physical minibatch per replica.
For example, \(K=32,B=64\) estimates about 18.0 GiB for learner-side
persistent state and activations, but that number is not a whole-session
admission proof. It was fit from a narrow set of matched H100 runs and
omits other phases. The correct workflow is to use it to choose a
candidate, compile a full lifetime envelope, then validate the warmed
steady-state high-water mark.

\subsection{18. H100 versus RTX 2080 Ti: not the same GPU at different
speed}\label{18-h100-versus-rtx-2080-ti-not-the-same-gpu-at-different-speed}

The local server uses RTX 2080 Ti hardware for development and CUDA
regression; the paid/cluster target is an NVIDIA H100 80 GB HBM3. The
relevant H100 run identified an SXM-class device with compute capability
9.0, 85,017,493,504 bytes visible to PyTorch, BF16 support, and a 700 W
power cap. The comparison below uses NVIDIA reference/Founders Edition
RTX 2080 Ti figures and H100 SXM specifications; installed board clocks
and software can differ.

{\def\LTcaptype{none} % do not increment counter
\begin{longtable}[]{@{}
  >{\raggedright\arraybackslash}p{0.16\linewidth}
  >{\raggedright\arraybackslash}p{0.14\linewidth}
  >{\raggedright\arraybackslash}p{0.22\linewidth}
  >{\raggedright\arraybackslash}p{0.35\linewidth}@{}}
\toprule\noalign{}
Property & RTX 2080 Ti & H100 SXM 80 GB & Consequence here \\
\midrule\noalign{}
\endhead
\bottomrule\noalign{}
\endlastfoot
Architecture & Turing TU102 & Hopper GH100 & Different kernels and
tuning rules \\
Compute capability & 7.5 (SM75) & 9.0 (SM90) & Extensions need correct
targets/caches \\
SM count & 68 & 132 enabled in SXM H100 & More parallel capacity to
feed \\
Device memory & 11 GB GDDR6 & 80 GB HBM3 & About \(7.3\times\) nominal
capacity \\
Peak memory bandwidth & 616 GB/s & 3.35 TB/s & About \(5.4\times\)
specification bandwidth \\
Peak non-Tensor FP32 & 14.2 TFLOP/s FE & 67 TFLOP/s & About
\(4.7\times\) specification peak \\
Low-precision Tensor math & FP16; 113.8 Tensor TFLOP/s FE & BF16/FP16;
989 dense or 1,979 sparse Tensor TFLOP/s & H100 adds BF16 and a much
larger matrix ceiling \\
BF16 Tensor Core support & No & Yes & Local BF16 production path cannot
run \\
Nominal board/TDP & 260 W FE & up to 700 W & Different thermal/power
environment \\
Memory technology & GDDR6 & HBM3 & Capacity and bandwidth change
geometry choices \\
\end{longtable}
}

H100\textquotesingle s published low-precision figures require careful
reading: NVIDIA lists 1,979 BF16/FP16 Tensor TFLOP/s with sparsity for
H100 SXM, implying 989 TFLOP/s for the corresponding dense
specification. Neither figure predicts this project\textquotesingle s
throughput. Poker routing, small GEMMs, attention shapes, launch gaps,
completion, targets, and optimizer work all sit outside a single dense
matrix peak.

\subsubsection{18.1 What the 2080 Ti is good
for}\label{181-what-the-2080-ti-is-good-for}

The local Turing GPU is valuable for:

\begin{itemize}
\tightlist
\item
  CUDA residency and no-CPU-fallback gates;
\item
  FP32 control-flow, storage, target, evaluator, and learner tests;
\item
  memory-safety checks at small geometry;
\item
  deterministic sampling and transaction/rollback tests;
\item
  SM75 extension compilation; and
\item
  cheap iteration on profiler annotations and benchmark harnesses.
\end{itemize}

Its smaller capacity can even expose memory waste earlier. An
implementation that only works because an H100 has 80 GB is not
necessarily well designed.

\subsubsection{18.2 What it cannot
establish}\label{182-what-it-cannot-establish}

The RTX 2080 Ti cannot establish:

\begin{itemize}
\tightlist
\item
  BF16 forward/backward correctness or performance;
\item
  SM90-native extension compatibility;
\item
  Hopper-specific Flash/SDPA or grouped-GEMM dispatch;
\item
  H100 Tensor Core efficiency;
\item
  the root/replica geometry admitted by 80 GB;
\item
  HBM3 behavior;
\item
  Hopper concurrency or launch/compute balance; or
\item
  the exact framework compatibility surface on the rental image.
\end{itemize}

This is why the local pilot deliberately used FP32, tiny roots, batch 8,
one gradient step, and a small RB completion count. Those are smoke
semantics, not H100 throughput estimates.

\subsubsection{18.3 Newer hardware can expose more
bugs}\label{183-newer-hardware-can-expose-more-bugs}

Moving to H100 activates paths that local testing never executes:

\begin{itemize}
\tightlist
\item
  BF16 storage and autocast;
\item
  much larger cumulative geometries;
\item
  SM90 compilation and kernel selection;
\item
  library implementations chosen only for newer hardware;
\item
  enough concurrency to expose race/order defects; and
\item
  longer or more expensive arms where missing progress telemetry
  matters.
\end{itemize}

It also makes host overhead more visible. If a fixed CPU launch sequence
takes the same time on both machines while each GPU kernel becomes
faster, the H100 spends a larger fraction of time waiting. Peak hardware
magnifies systems inefficiency.

\subsection{19. What the H100 debugging history actually taught
us}\label{19-what-the-h100-debugging-history-actually-taught-us}

The repository contains two related H100 episodes: an initial late-July
pilot and a broader August 2 RunPod selection. The following is a
compact chronology, not a claim that every historical branch equals the
current checkout.

\subsubsection{19.1 Baseline rewrite and initial H100
pilot}\label{191-baseline-rewrite-and-initial-h100-pilot}

By July 28, the queued accelerator-only engine, fixed storage,
transactional rollout, evaluator integration, and layered targets had
replaced earlier execution surfaces. H100 pilot tooling and Normal/RB
configs followed on July 29. A RunPod pilot proved that one-replica HU
32 bb Normal and RB could execute in BF16 with 512 roots, logical batch
256, and 116 learner steps. The audit reported about 42.13 Normal and
41.12 RB player nodes/s, with peak reserved memory no greater than 16.38
GiB for those successful arms.

That success did not mean the launch path was clean. It produced four
concrete blockers/lessons.

\subsubsection{19.2 Failure 1: FP32 model input routed into BF16
storage}\label{192-failure-1-fp32-model-input-routed-into-bf16-storage}

\textbf{Symptom.} A BF16 smoke failed when strict \texttt{index\_copy\_}
attempted to write an FP32 \texttt{model\_input} into a BF16 family
arena.

\textbf{Cause.} Upstream model-input construction intentionally produced
FP32, while the H100 execution manifest declared BF16 model storage. The
route boundary had not encoded the allowed conversion.

\textbf{Repair.} Cast only floating-point \texttt{model\_input} at that
declared boundary. Keep every other dtype mismatch fatal, especially
integer addresses, owners, masks, and metadata.

\textbf{General lesson.} Mixed precision needs typed interface
boundaries. ``Make PyTorch cast whatever is convenient'' is unsafe;
``never cast'' is also wrong when storage precision is an explicit
design choice.

The validated RunPod repair existed first as \texttt{2ebfae7} and was
ported into the main line by \texttt{dcb120a} on July 31 with focused
tests.

\subsubsection{19.3 Failure 2: a cumulative folded-range owner sized as
one
pack}\label{193-failure-2-a-cumulative-folded-range-owner-sized-as-one-pack}

\textbf{Symptom.} With 512 roots and a scheduler pack of 84, the
folded-range slab overflowed deterministically when the fifth pack began
at logical root 336.

\textbf{Cause.} Configuration used roughly
\texttt{min(family\_pack,\ roots)} rows for data that survived across
the complete rollout. That formula described simultaneous producer work,
not cumulative retained ownership.

\textbf{Repair.} For the pilot, size the per-partition folded capacity
by the full-root lifetime; the generalized bound accounts for every
possible folder. Later manifest compilation derived such capacities
centrally rather than requiring duplicated YAML arithmetic.

\textbf{General lesson.} Capacity is a theorem about lifetime. The error
was not ``H100 needs more memory''; it was ``the formula bounded the
wrong object.'' Adding a generous constant would have hidden the
misunderstanding.

The validated pilot repair was \texttt{d9db16c}; the main-line port
again landed in \texttt{dcb120a}.

\subsubsection{19.4 Failure 3: K replicas performed the learner
twice}\label{194-failure-3-k-replicas-performed-the-learner-twice}

\textbf{Symptom.} Normal K=1 took 39.69 seconds while K=2 took 158.82
seconds; RB K=1 took 41.52 seconds while K=2 again took 158.82 seconds.
Doubling work caused approximately \(4\times\) time, so aggregate
throughput fell to about half.

\textbf{Cause.} The replica/distributed path first ran a complete
validate-only forward/loss/backward for every lane, cleared it, and then
repeated the same work for committed updates. Lanes were also processed
serially. K=4 reached high sampled GPU utilization but timed out; it was
doing abundant redundant work, not demonstrating efficient parallel
learning.

\textbf{Repair.} Commit \texttt{9229734} removed the duplicated learner
preflight and introduced one transactional learner pass per lane with
rollback/failure agreement over model, optimizer, RNG, replica version,
and run state.

\textbf{General lesson.} Utilization and correctness machinery must be
included in the work ledger. A GPU at 95\% utilization can be
productively wrong, or merely busy computing the same gradients twice.

\subsubsection{19.5 Failure 4: long jobs produced insufficient
evidence}\label{195-failure-4-long-jobs-produced-insufficient-evidence}

\textbf{Symptom.} Timed-out arms could leave empty logs and no partial
benchmark report. Aggregate timing did not explain whether rollout,
completion, targets, learning, or synchronization was responsible.

\textbf{Repair.} The benchmark added atomic progress reports at
warmup/repeat boundaries. Later H100 work used detached launches,
per-arm exit files, metrics, GPU monitoring, and harvested artifacts.

\textbf{General lesson.} Observability is production functionality. A
benchmark that cannot explain or resume a failure wastes both hardware
rental time and the opportunity to improve the model.

\subsubsection{19.6 Correctness found during performance
preparation}\label{196-correctness-found-during-performance-preparation}

The H100 preparation also found a statistical issue. An automatic
cost-aware sequential allocation favored cheap preflop rows (roughly
239/17 or 238/18 pre/post in representative retained data) but weighted
the learner as though the allocation were the target distribution. This
changed the effective decision objective.

Commit \texttt{ef22ee2} changed default fixed-step sequential training
to sample with replacement from the pooled retained decision extent,
grouping by family only for physical model execution. Explicit
cost-stratified allocation remains possible with normalized target-mass
weights, and physical microbatch slices retain exact logical-batch
weight.

\textbf{General lesson.} Performance-oriented stratification is
importance sampling whether or not the code calls it that. Proposal
probabilities and target mass must be reconciled mathematically before
comparing throughput.

\subsubsection{19.7 August 2 environment compatibility
surprises}\label{197-august-2-environment-compatibility-surprises}

The larger RunPod experiment started from \texttt{ed4a73e} on an H100 80
GB HBM3 with driver 580.126.09, PyTorch 2.4.1+cu124, CUDA
runtime/toolkit 12.4, SM90, BF16, and MIG disabled. Preflight found two
more operational defects:

\begin{enumerate}
\def\labelenumi{\arabic{enumi}.}
\tightlist
\item
  Collecting framework and evaluator tests in one pytest process caused
  top-level \texttt{tests} package shadowing. Running the suites in
  separate processes fixed the invocation without source changes.
\item
  PyTorch 2.4.1 CUDA did not support stable \texttt{argsort} directly on
  boolean keys. Casting the sort key to \texttt{int32} preserved
  \texttt{False=0}, \texttt{True=1}, stable row order, and scientific
  semantics. This became experiment commit \texttt{3e563319}.
\end{enumerate}

A reporting bug used nonexistent \texttt{IterationResult.steps} instead
of \texttt{optimizer\_steps}, so the explicit gradient telemetry block
was absent despite successful finite updates. The experiment fixed that
in \texttt{8e305ca}. Neither experiment commit was present at the
documented starting main-line revision; the handoff retained a bundle
and patches rather than claiming they had been pushed.

\textbf{General lesson.} Framework version is part of the executable
contract. A valid operation on one CUDA/PyTorch stack may have a
dtype-specific hole on another. Preflight should exercise real dtypes
and operators before launching the expensive matrix.

\subsection{20. What the August H100 experiment
measured}\label{20-what-the-august-h100-experiment-measured}

The August 2 experiment was explicitly a fresh-initialization execution
benchmark. It did not compare trained policy quality, exploitability,
convergence, or estimator equivalence. Every launched experimental arm
exited successfully; the final Normal and RB reliability runs each
completed three iterations, 116 finite-gradient optimizer commits per
iteration, checkpoints, and clean shutdown.

The selected HU 32 bb scientific/execution baseline was:

{\def\LTcaptype{none} % do not increment counter
\begin{longtable}[]{@{}lr@{}}
\toprule\noalign{}
Setting & Selected value \\
\midrule\noalign{}
\endhead
\bottomrule\noalign{}
\endlastfoot
replicas \(K\) & 1 \\
roots & 4,096 \\
rollout/completion/learner packs & 256 / 512 / 256 \\
logical learner batch & 256 \\
learner steps & 116 \\
precision & BF16 autocast with FP32 training state \\
attention & SDPA \\
reference cache & auto \\
RB completion samples & 1,024 \\
\end{longtable}
}

The pilot model used \(d_{\rm model}=128\), four heads, four preflop
layers, six postflop layers, and two raise components.

\subsubsection{20.1 The decisive result was root parallelism, not serial
replicas}\label{201-the-decisive-result-was-root-parallelism-not-serial-replicas}

Fair cache-off K comparisons at 512 roots retained K=1:

{\def\LTcaptype{none} % do not increment counter
\begin{longtable}[]{@{}lrrrl@{}}
\toprule\noalign{}
Mode & K=1 & K=2 & Aggregate K=2 gain & Decision \\
\midrule\noalign{}
\endhead
\bottomrule\noalign{}
\endlastfoot
Normal & 13.220 nodes/s & 13.345 nodes/s & 0.95\% & keep K=1 \\
RB & 12.837 nodes/s & 13.235 nodes/s & 3.10\% & keep K=1 \\
\end{longtable}
}

Those gains were below the predeclared 10\% threshold. Short K=4 probes
at 64 roots per replica passed stability/memory gates but were not fair
throughput comparisons; K=8 was not run.

Increasing roots was dramatically more valuable:

{\def\LTcaptype{none} % do not increment counter
\begin{longtable}[]{@{}lrrr@{}}
\toprule\noalign{}
Mode & 512 roots & 2,048 roots & 4,096 roots \\
\midrule\noalign{}
\endhead
\bottomrule\noalign{}
\endlastfoot
Normal & 23.662 nodes/s & 83.442 nodes/s & 158.121 nodes/s \\
RB & 23.069 nodes/s & 85.160 nodes/s & 153.293 nodes/s \\
\end{longtable}
}

The arithmetic makes the Pareto result concrete. Multiplying roots by
eight from 512 to 4,096 multiplied Normal throughput by about \(6.68\)
and RB throughput by about \(6.64\). That is roughly 84\% and 83\% of
ideal linear scaling over this interval. In contrast, doubling serial
replicas improved aggregate throughput by only 0.95\% and 3.10\%. The
root sweep did not make each root cheaper in every abstract sense; it
gave the H100 enough independent rows to amortize fixed costs and occupy
more of the machine.

This does not prove replica batching is a bad idea. The tested K path
remained largely serial at the learner level, so it mostly replicated
work and state. It proves that, for the implementation actually
measured, using H100 memory to expose more root-level work was much
higher ROI than adding serial replicas.

\subsubsection{20.2 Steady-state memory was much higher than
one-iteration
memory}\label{202-steady-state-memory-was-much-higher-than-one-iteration-memory}

The selected single-iteration arms reported first-iteration peak
reserved memory of 26.498 GiB Normal and 34.736 GiB RB. Three-iteration
reliability runs, which matured the automatic reference cache, reached
conservative cumulative peaks of 50.271 GiB and 58.619 GiB.

Normal cumulative peak plateaued by iteration 3. RB changed by only
about 50.3 MB (0.08\%) from iteration 2 to 3 while end-reserved memory
decreased, consistent with stable workload/allocator variation rather
than sustained growth.

This is a textbook memory-lifetime lesson: a one-step peak was not the
session peak because a later persistent cache changed the live set.

The measured H100 exposed about 79.18 GiB to PyTorch. The 58.619 GiB RB
peak therefore left about 20.6 GiB between the recorded high-water mark
and visible capacity. That difference is not an invitation to allocate
20.6 GiB: the design still reserves 10 GiB for safety, and future
shapes, library workspaces, fragmentation, compiled graphs, or cache
maturation can consume part of the rest. Admission must compare the full
phase envelope with the configured reserve, not merely check that the
previous run did not OOM.

\subsubsection{20.3 Optional knobs did not justify scientific
changes}\label{203-optional-knobs-did-not-justify-scientific-changes}

Reducing learner steps from 116 to 58 shortened wall time and increased
rollout nodes/s, but reduced learner presentations/s. That is a schedule
change, so the experiment retained 116 without learning-quality
evidence.

At logical learner batch 256, one replica performs

\[
116\cdot256=29,696
\]

learner presentations per iteration, or 89,088 across the
three-iteration reliability run. A 58-step arm performs exactly half
that learning work. Its shorter wall time is useful capacity-planning
evidence, but calling it an executor speedup would compare different
algorithms.

Reducing RB completion samples from 1,024 to 512 improved measured
nodes/s by only about 0.8\%. Because sample count changes estimator
quality and no variance or exploitability equivalence was established,
the experiment retained 1,024.

\subsubsection{20.4 The final measured
baseline}\label{204-the-final-measured-baseline}

Across the three-iteration steady-state confirmations, the handoff
reported:

\begin{itemize}
\tightlist
\item
  Normal: 169.563 player nodes/s and 459.778 learner presentations/s at
  a 50.271 GiB measured steady-state peak;
\item
  RB: 174.111 player nodes/s and 456.481 learner presentations/s at a
  58.619 GiB measured steady-state peak.
\end{itemize}

These are execution baselines for that exact fresh-initialization
environment, not general H100 claims. The formal Flash-attention
dispatch proof remained open, and a PyTorch \texttt{vmap}
batching-fallback warning appeared for efficient SDPA in targeted tests.

\subsection{21. A Pareto-ranked experiment and optimization
backlog}\label{21-a-pareto-ranked-experiment-and-optimization-backlog}

The following is a proposal, not a description of already implemented
production behavior. Its order is chosen to maximize information and
reusable value before specialized engineering.

\subsubsection{Tier 0 --- Reconcile and reproduce the exact
target}\label{tier-0--reconcile-and-reproduce-the-exact-target}

\begin{enumerate}
\def\labelenumi{\arabic{enumi}.}
\tightlist
\item
  \textbf{Reconcile experiment-only fixes.} Decide whether the
  boolean-sort and gradient-telemetry commits from the August handoff
  belong on the intended branch. Run the full CUDA gates after any port;
  do not treat a patch archive as current source.
\item
  \textbf{Reproduce the selected 4,096-root configs on the target Quest
  H100.} Record exact SKU, driver, toolkit, PyTorch/cuDNN, evaluator
  revision, native SM90 extension identity, clocks, MIG state, and
  contention.
\item
  \textbf{Prove SDPA dispatch.} Force or inspect Flash forward and
  backward on real preflop/postflop shapes. If a shape falls back,
  record why and its runtime coverage before changing the model.
\item
  \textbf{Validate checkpoint/resume and trained initialization.} The
  August run was fresh-initialization execution evidence. A scientific
  run needs preserved model provenance and resume identity.
\item
  \textbf{Add phase-complete progress.} Warmup/repeat reports are
  useful; longer production runs also benefit from phase and
  learner-step heartbeats without introducing hot-path synchronization.
\end{enumerate}

These steps are high ROI because every later speed comparison depends on
them.

\subsubsection{Tier 1 --- Improve the measured K=1 critical
path}\label{tier-1--improve-the-measured-k1-critical-path}

\begin{enumerate}
\def\labelenumi{\arabic{enumi}.}
\tightlist
\item
  Capture an unprofiled phase breakdown and an Nsight Systems trace for
  a fixed learner replay and one full iteration.
\item
  Separate learner core time (staged inputs through optimizer) from
  gathering, rollout, completion, targets, reference cache, and
  checkpoint work.
\item
  Record every dominant GEMM/attention shape and exact physical/useful
  token count. Compare kernels with isolated same-shape ceilings.
\item
  Sweep a narrow neighborhood around 4,096 roots and the 256/512/256
  packs, subject to the 10 GiB reserve and steady-state cache peak. The
  earlier sweep found a steeply rising curve; it did not necessarily
  locate the plateau.
\item
  Compile stable learner regions and compare eager, compiled, and---only
  for a bounded static shape set---CUDA Graph replay. Track compilation
  count, graph pool memory, and numerical drift.
\item
  Vectorize finite checks, norm/clipping, gradient clearing, and Adam
  operations where profiling shows per-parameter or per-lane launch
  amplification.
\end{enumerate}

The retention criterion should be end-to-end or full-replay gain, not an
operator microbenchmark alone.

\subsubsection{Tier 2 --- Make replicas visible to the
GPU}\label{tier-2--make-replicas-visible-to-the-gpu}

The existing K=2 result does not answer whether a truly batched replica
learner could saturate H100 better. Build one immutable replay artifact
so candidates receive identical per-lane inputs, targets, sampling
order, initial parameters, moments, and RNG.

Compare these implementations:

\begin{enumerate}
\def\labelenumi{\arabic{enumi}.}
\tightlist
\item
  \textbf{Serial oracle:} current one-lane-at-a-time learner.
\item
  \textbf{CUDA stream waves:} unchanged per-lane model on multiple
  streams; cheapest test of whether independent kernels leave
  complementary resources idle.
\item
  \textbf{Functional \texttt{vmap} plus outer backward:} one canonical
  model with stacked state; reject batching fallbacks and loss of Flash
  backward.
\item
  \textbf{Explicit replica-native BMM:} \texttt{{[}K,...{]}} weights,
  batched linear primitives, and replica folded into the SDPA batch.
\item
  \textbf{Grouped GEMM for tails:} only if measured padding waste is
  material and the target software/SM90 kernel is stable.
\end{enumerate}

The scheduler should be \textbf{step-major}, not lane-major:

\begin{verbatim}
for logical step s:
    prepare the next batch for every active lane
    execute the lane cohort
    validate all lanes/ranks
    apply independent per-lane transitions
\end{verbatim}

Within each lane, step \(s\) must commit before step \(s+1\). Across
lanes, work can be concurrent. Loss normalization, clipping, moments,
RNG, and rollback remain lane-local. A summed cohort loss is acceptable
for one outer backward only if it is a sum of independent lane losses
and gradients land in disjoint stacked parameter rows; averaging across
\(K\) would alter each transition.

Suggested gates before production retention:

\begin{itemize}
\tightlist
\item
  K=1 regression no worse than 3\%;
\item
  K=2 aggregate learner gain at least 25\%;
\item
  K=4 gain at least 75\%;
\item
  at least 10\% over the next-best correct new executor;
\item
  exact logical samples and per-lane isolation;
\item
  FP32 one-step oracle comparison and predeclared BF16 error bounds;
\item
  complete 116-step drift, checkpoint, and failure-injection gates;
\item
  Flash forward/backward retained; and
\item
  no memory growth with the configured reserve.
\end{itemize}

These are proposed decision thresholds, not promises that the hardware
will meet them.

\subsubsection{Tier 3 --- Reduce model-side wasted work without mixing
the
experiment}\label{tier-3--reduce-model-side-wasted-work-without-mixing-the-experiment}

Potentially valuable ideas include:

\begin{itemize}
\tightlist
\item
  compact board-impossible postflop hand rows before expensive
  attention;
\item
  evaluate whether the full all-hands readout is required for every
  Normal presentation;
\item
  reduce avoidable pre/post padding through measured buckets;
\item
  remove redundant incidence/folded features in a heads-up-specific
  model; and
\item
  revisit long-sequence attention structure.
\end{itemize}

Some of these alter architecture or the domain of outputs. They must be
tested as separate scientific factors with learning-quality evidence. A
faster heads-up-lite model cannot establish that an executor is faster,
and it cannot silently replace the common Normal/RB comparison model.

\subsubsection{Tier 4 --- Custom kernels after
profiling}\label{tier-4--custom-kernels-after-profiling}

Candidate custom work already suggested by design audits includes:

\begin{itemize}
\tightlist
\item
  fused posterior/profile sampling without a large dense attempt
  envelope;
\item
  direct selected-action policy computation rather than materializing
  every legal-action probability;
\item
  fused transition/folded-statistic construction;
\item
  explicit replica linear primitives where library batching is poor; and
\item
  persistent kernels only for a demonstrated dominant fixed shape.
\end{itemize}

For each candidate calculate Amdahl coverage from a real trace. Require
a credible full-step upside before implementation and a material
measured gain after integration. Keep one ordinary correctness oracle
and delete losing production backends after the experiment; a permanent
forest of switches is a systems cost.

\subsection{22. A concrete experiment
matrix}\label{22-a-concrete-experiment-matrix}

The safest way to spend scarce H100 time is a staged, resumable matrix.

\subsubsection{Stage A: local RTX 2080 Ti
gates}\label{stage-a-local-rtx-2080-ti-gates}

\begin{itemize}
\tightlist
\item
  Run FP32 one-step and 116-step replays for every candidate available
  on SM75.
\item
  Test exact sampling, lane isolation, rollback, checkpoint/resume, and
  memory reconciliation.
\item
  Compile extensions into an SM75-specific cache.
\item
  Runtime-gate SM90-only grouped kernels rather than pretending the
  local test covered them.
\end{itemize}

\subsubsection{Stage B: five-minute H100
preflight}\label{stage-b-five-minute-h100-preflight}

\begin{itemize}
\tightlist
\item
  Confirm an exclusive non-MIG H100, identity, power/clocks, and clean
  memory.
\item
  Compile/load every native extension from an SM90-specific cache.
\item
  Run tiny BF16 Normal and RB forward/backward/completion.
\item
  Verify real Flash attention forward/backward dispatch.
\item
  Verify profiler permissions and every output/progress path.
\end{itemize}

\subsubsection{Stage C: correctness and
dispatch}\label{stage-c-correctness-and-dispatch}

\begin{itemize}
\tightlist
\item
  K=1 serial oracle for full preflop and postflop widths.
\item
  K=2 one-step and complete-replay tests for each executor.
\item
  Isolated forward, input-gradient, and weight-gradient benchmarks for
  dominant linear shapes.
\item
  Explicit useful versus padded work accounting.
\end{itemize}

\subsubsection{Stage D: broad unprofiled
sweep}\label{stage-d-broad-unprofiled-sweep}

\begin{itemize}
\tightlist
\item
  Sweep admitted \(K\in\{1,2,4,8\}\) and physical learner cohorts around
  16/32/64/128 rows where meaningful.
\item
  Sweep fixed total cohort work \(KB\) to separate ``more replicas''
  from ``more model work.''
\item
  Include preflop, postflop, mixed learner replay, and complete
  iteration.
\item
  Record steady-state timing and memory after caches/graphs mature.
\end{itemize}

\subsubsection{Stage E: profile only
finalists}\label{stage-e-profile-only-finalists}

\begin{itemize}
\tightlist
\item
  Capture bounded Nsight Systems traces with phase annotations.
\item
  Profile representative GEMM/BMM, Flash forward/backward, incidence,
  reduction, and optimizer kernels with Nsight Compute.
\item
  Compare achieved performance to isolated same-shape ceilings and the
  relevant rooflines.
\end{itemize}

\subsubsection{Stage F: reliability and scientific
handoff}\label{stage-f-reliability-and-scientific-handoff}

\begin{itemize}
\tightlist
\item
  Alternate at least several complete Normal/RB repeats.
\item
  Run multiple iterations to expose persistent caches and allocator
  growth.
\item
  Verify checkpoints, resume, exact work ledgers, finite metrics, and
  clean teardown.
\item
  Report speed, memory, correctness, and scientific limitations
  together.
\end{itemize}

Stop a candidate immediately when it changes logical work, crosses a
numerical gate, couples replica state, loses required attention
dispatch, exceeds memory, or cannot explain its speedup. A failed
\texttt{vmap} experiment does not disprove explicit BMM; a failed BMM
experiment does not disprove stream waves. It only answers the specific
implementation tested.

\subsection{23. Closing perspective}\label{23-closing-perspective}

The main lesson of the project is not that H100 is ``faster'' than RTX
2080 Ti. It is that a larger accelerator changes the active constraints.

The 2080 Ti is an excellent correctness laboratory for common CUDA
behavior, but it cannot execute the H100 BF16 and SM90 contract. The
H100\textquotesingle s extra memory made 4,096-root runs possible and
those roots delivered far more aggregate value than serial replica
scaling. At the same time, H100-specific execution exposed a missing
precision cast, a lifetime-bound capacity error, framework operator
compatibility, insufficient progress telemetry, and redundant learner
work. None of those problems would be solved by quoting Tensor TFLOPS.

The project\textquotesingle s highest-return systems choices are
correspondingly plain:

\begin{itemize}
\tightlist
\item
  define and count correct statistical work;
\item
  derive one immutable capacity plan from one config;
\item
  keep production state accelerator-resident in typed, lifetime-aware
  storage;
\item
  use device transactions and logical RNG so scheduling can change
  safely;
\item
  call optimized attention/evaluator primitives and prove their
  dispatch;
\item
  increase the cheapest honest parallel axis until throughput plateaus;
\item
  measure steady state, including caches and rollback state; and
\item
  specialize kernels only where a real trace and
  Amdahl\textquotesingle s law justify them.
\end{itemize}

That is the Pareto view of GPU engineering: understand the bottleneck at
the scale of the whole system, make one falsifiable change, retain it
only when the useful workload becomes materially better, and let
evidence---not the novelty of the technique---choose the design.

\begin{center}\rule{0.5\linewidth}{0.5pt}\end{center}

\subsection{Glossary}\label{glossary}

\textbf{Arithmetic intensity:} Operations performed per byte moved from
a chosen memory level.

\textbf{Autocast:} Framework policy that selects lower precision for
eligible operations while leaving other operations in higher precision.

\textbf{BF16 / bfloat16:} A 16-bit floating format with an FP32-like
exponent range and reduced significand precision.

\textbf{Block / CTA:} A cooperative group of CUDA threads scheduled
together on one SM. CTA means cooperative thread array and is commonly
used as a synonym for thread block.

\textbf{CUDA Graph:} An instantiated device-operation dependency graph
designed for lower-overhead replay of a stable schedule.

\textbf{Device:} The GPU in CUDA terminology; the CPU is the host.

\textbf{GEMM / BMM:} General matrix multiplication / batched matrix
multiplication.

\textbf{HBM:} High-bandwidth memory, the device DRAM technology used by
the H100 SXM here.

\textbf{Kernel:} A function launched for parallel execution on the GPU.

\textbf{Logical batch:} The examples defining one statistical update,
independent of physical slicing.

\textbf{Microbatch / pack:} A bounded physical subset executed
concurrently to control memory or shape.

\textbf{Occupancy:} Resident warps relative to an SM\textquotesingle s
architectural maximum, constrained by block resources. It is a
diagnostic, not the optimization objective.

\textbf{Rao-\/-Blackwellization (RB):} Replacing some sampled
uncertainty with a conditional expectation, typically reducing estimator
variance at additional compute/memory cost.

\textbf{Replica:} One independent model, optimizer, RNG, and run state.
Replica batching must not share weights or reduce gradients across
replicas.

\textbf{SDPA:} Scaled dot-product attention; a high-level operator that
may dispatch to Flash, memory-efficient, or math implementations
depending on its inputs and stack.

\textbf{SIMT / warp:} CUDA\textquotesingle s single-instruction,
multiple-thread execution model; a warp contains 32 threads.

\textbf{SM / SM75 / SM90:} Streaming multiprocessor; SM75 and SM90
informally denote compute capabilities 7.5 (Turing) and 9.0 (Hopper).

\textbf{SoA:} Structure of arrays, a layout placing each field in its
own dense array.

\textbf{Tensor Core:} Specialized hardware for supported matrix
multiply-accumulate operations.

\subsection{References and repository
evidence}\label{references-and-repository-evidence}

Generic primary references:

\begin{itemize}
\tightlist
\item
  \href{https://docs.nvidia.com/cuda/cuda-programming-guide/contents.html}{NVIDIA
  CUDA Programming Guide}
\item
  \href{https://docs.nvidia.com/cuda/cuda-c-best-practices-guide/index.html}{NVIDIA
  CUDA C++ Best Practices Guide}
\item
  \href{https://docs.nvidia.com/cuda/hopper-tuning-guide/index.html}{NVIDIA
  Hopper Tuning Guide}
\item
  \href{https://docs.nvidia.com/cuda/turing-tuning-guide/index.html}{NVIDIA
  Turing Tuning Guide}
\item
  \href{https://developer.nvidia.com/blog/nvidia-hopper-architecture-in-depth/}{NVIDIA
  Hopper architecture in depth}
\item
  \href{https://images.nvidia.com/aem-dam/Solutions/design-visualization/technologies/turing-architecture/NVIDIA-Turing-Architecture-Whitepaper.pdf}{NVIDIA
  Turing architecture whitepaper}
\item
  \href{https://www.nvidia.com/en-us/data-center/h100/}{NVIDIA H100
  specifications}
\item
  \href{https://www.nvidia.com/content/nvidiaGDC/zz/en_ZZ/geforce/graphics-cards/rtx-2080-ti.html}{NVIDIA
  GeForce RTX 2080 Ti specifications}
\item
  \href{https://docs.pytorch.org/docs/stable/generated/torch.nn.functional.scaled_dot_product_attention}{PyTorch
  scaled dot-product attention}
\item
  \href{https://docs.pytorch.org/docs/stable/generated/torch.compile.html}{PyTorch
  \texttt{torch.compile}}
\end{itemize}

Project evidence used for the case study:

\begin{itemize}
\tightlist
\item
  \texttt{README.md}
\item
  \texttt{docs/queued/core\_pipeline.md}
\item
  \texttt{docs/queued/storage\_and\_lifetimes.md}
\item
  \texttt{docs/queued/topology\_and\_targets.md}
\item
  \texttt{docs/queued/h100\_pilot.md}
\item
  \texttt{docs/audit/queued\_issues/quest\_h100\_launch\_blockers\_20260729.md}
\item
  \texttt{docs/audit/h100\_normal\_rb\_execution\_changes\_20260802.md}
\item
  \texttt{runs/experiments/nlhe\_hu\_32bb\_h100\_pilot/run\_state*.json}
\item
  \texttt{runs/nlhe\_hu\_32bb\_runpod\_h100\_20260802\_handoff/\{RESULTS,DECISIONS,AUDIT,HANDOFF\}.md}
\item
  Git history through \texttt{ed4a73e}, plus the two experiment-only
  handoff commits \texttt{3e563319} and \texttt{8e305ca} recorded in
  that bundle
\end{itemize}

The untracked \texttt{docs/audit/h100\_idea\_*.md} and
\texttt{docs/audit/h100\_plan.md} files were consulted as proposal
material only. This chapter does not treat their suggested executors,
thresholds, or model variants as implemented production features.

\end{document}
